% Options for packages loaded elsewhere
\PassOptionsToPackage{unicode}{hyperref}
\PassOptionsToPackage{hyphens}{url}
%
\documentclass[
]{book}
\usepackage{amsmath,amssymb}
\usepackage{iftex}
\ifPDFTeX
  \usepackage[T1]{fontenc}
  \usepackage[utf8]{inputenc}
  \usepackage{textcomp} % provide euro and other symbols
\else % if luatex or xetex
  \usepackage{unicode-math} % this also loads fontspec
  \defaultfontfeatures{Scale=MatchLowercase}
  \defaultfontfeatures[\rmfamily]{Ligatures=TeX,Scale=1}
\fi
\usepackage{lmodern}
\ifPDFTeX\else
  % xetex/luatex font selection
\fi
% Use upquote if available, for straight quotes in verbatim environments
\IfFileExists{upquote.sty}{\usepackage{upquote}}{}
\IfFileExists{microtype.sty}{% use microtype if available
  \usepackage[]{microtype}
  \UseMicrotypeSet[protrusion]{basicmath} % disable protrusion for tt fonts
}{}
\makeatletter
\@ifundefined{KOMAClassName}{% if non-KOMA class
  \IfFileExists{parskip.sty}{%
    \usepackage{parskip}
  }{% else
    \setlength{\parindent}{0pt}
    \setlength{\parskip}{6pt plus 2pt minus 1pt}}
}{% if KOMA class
  \KOMAoptions{parskip=half}}
\makeatother
\usepackage{xcolor}
\usepackage{color}
\usepackage{fancyvrb}
\newcommand{\VerbBar}{|}
\newcommand{\VERB}{\Verb[commandchars=\\\{\}]}
\DefineVerbatimEnvironment{Highlighting}{Verbatim}{commandchars=\\\{\}}
% Add ',fontsize=\small' for more characters per line
\usepackage{framed}
\definecolor{shadecolor}{RGB}{248,248,248}
\newenvironment{Shaded}{\begin{snugshade}}{\end{snugshade}}
\newcommand{\AlertTok}[1]{\textcolor[rgb]{0.94,0.16,0.16}{#1}}
\newcommand{\AnnotationTok}[1]{\textcolor[rgb]{0.56,0.35,0.01}{\textbf{\textit{#1}}}}
\newcommand{\AttributeTok}[1]{\textcolor[rgb]{0.13,0.29,0.53}{#1}}
\newcommand{\BaseNTok}[1]{\textcolor[rgb]{0.00,0.00,0.81}{#1}}
\newcommand{\BuiltInTok}[1]{#1}
\newcommand{\CharTok}[1]{\textcolor[rgb]{0.31,0.60,0.02}{#1}}
\newcommand{\CommentTok}[1]{\textcolor[rgb]{0.56,0.35,0.01}{\textit{#1}}}
\newcommand{\CommentVarTok}[1]{\textcolor[rgb]{0.56,0.35,0.01}{\textbf{\textit{#1}}}}
\newcommand{\ConstantTok}[1]{\textcolor[rgb]{0.56,0.35,0.01}{#1}}
\newcommand{\ControlFlowTok}[1]{\textcolor[rgb]{0.13,0.29,0.53}{\textbf{#1}}}
\newcommand{\DataTypeTok}[1]{\textcolor[rgb]{0.13,0.29,0.53}{#1}}
\newcommand{\DecValTok}[1]{\textcolor[rgb]{0.00,0.00,0.81}{#1}}
\newcommand{\DocumentationTok}[1]{\textcolor[rgb]{0.56,0.35,0.01}{\textbf{\textit{#1}}}}
\newcommand{\ErrorTok}[1]{\textcolor[rgb]{0.64,0.00,0.00}{\textbf{#1}}}
\newcommand{\ExtensionTok}[1]{#1}
\newcommand{\FloatTok}[1]{\textcolor[rgb]{0.00,0.00,0.81}{#1}}
\newcommand{\FunctionTok}[1]{\textcolor[rgb]{0.13,0.29,0.53}{\textbf{#1}}}
\newcommand{\ImportTok}[1]{#1}
\newcommand{\InformationTok}[1]{\textcolor[rgb]{0.56,0.35,0.01}{\textbf{\textit{#1}}}}
\newcommand{\KeywordTok}[1]{\textcolor[rgb]{0.13,0.29,0.53}{\textbf{#1}}}
\newcommand{\NormalTok}[1]{#1}
\newcommand{\OperatorTok}[1]{\textcolor[rgb]{0.81,0.36,0.00}{\textbf{#1}}}
\newcommand{\OtherTok}[1]{\textcolor[rgb]{0.56,0.35,0.01}{#1}}
\newcommand{\PreprocessorTok}[1]{\textcolor[rgb]{0.56,0.35,0.01}{\textit{#1}}}
\newcommand{\RegionMarkerTok}[1]{#1}
\newcommand{\SpecialCharTok}[1]{\textcolor[rgb]{0.81,0.36,0.00}{\textbf{#1}}}
\newcommand{\SpecialStringTok}[1]{\textcolor[rgb]{0.31,0.60,0.02}{#1}}
\newcommand{\StringTok}[1]{\textcolor[rgb]{0.31,0.60,0.02}{#1}}
\newcommand{\VariableTok}[1]{\textcolor[rgb]{0.00,0.00,0.00}{#1}}
\newcommand{\VerbatimStringTok}[1]{\textcolor[rgb]{0.31,0.60,0.02}{#1}}
\newcommand{\WarningTok}[1]{\textcolor[rgb]{0.56,0.35,0.01}{\textbf{\textit{#1}}}}
\usepackage{longtable,booktabs,array}
\usepackage{calc} % for calculating minipage widths
% Correct order of tables after \paragraph or \subparagraph
\usepackage{etoolbox}
\makeatletter
\patchcmd\longtable{\par}{\if@noskipsec\mbox{}\fi\par}{}{}
\makeatother
% Allow footnotes in longtable head/foot
\IfFileExists{footnotehyper.sty}{\usepackage{footnotehyper}}{\usepackage{footnote}}
\makesavenoteenv{longtable}
\usepackage{graphicx}
\makeatletter
\def\maxwidth{\ifdim\Gin@nat@width>\linewidth\linewidth\else\Gin@nat@width\fi}
\def\maxheight{\ifdim\Gin@nat@height>\textheight\textheight\else\Gin@nat@height\fi}
\makeatother
% Scale images if necessary, so that they will not overflow the page
% margins by default, and it is still possible to overwrite the defaults
% using explicit options in \includegraphics[width, height, ...]{}
\setkeys{Gin}{width=\maxwidth,height=\maxheight,keepaspectratio}
% Set default figure placement to htbp
\makeatletter
\def\fps@figure{htbp}
\makeatother
\setlength{\emergencystretch}{3em} % prevent overfull lines
\providecommand{\tightlist}{%
  \setlength{\itemsep}{0pt}\setlength{\parskip}{0pt}}
\setcounter{secnumdepth}{5}
\usepackage{booktabs}
\ifLuaTeX
  \usepackage{selnolig}  % disable illegal ligatures
\fi
\usepackage[]{natbib}
\bibliographystyle{apalike}
\usepackage{bookmark}
\IfFileExists{xurl.sty}{\usepackage{xurl}}{} % add URL line breaks if available
\urlstyle{same}
\hypersetup{
  pdftitle={Conceptos de Superficies},
  pdfauthor={John Jairo Estrada Álvarez},
  hidelinks,
  pdfcreator={LaTeX via pandoc}}

\title{Conceptos de Superficies}
\author{John Jairo Estrada Álvarez}
\date{2026-04-26}

\begin{document}
\maketitle

{
\setcounter{tocdepth}{1}
\tableofcontents
}
\chapter{Vectores y Geometría Vectorial}\label{vectores}

🎧 Podcast del capítulo

Tu navegador no soporta el elemento de audio.

El estudio de las superficies en el espacio tridimensional descansa sobre los fundamentos del álgebra vectorial. En este capítulo se presentan los conceptos esenciales de vectores, sus operaciones y propiedades geométricas.

\section{Vectores libres y vectores coordenados}\label{vectores-libres-y-vectores-coordenados}

\subsection{Definición de vector}\label{definiciuxf3n-de-vector}

Un \textbf{vector} es un objeto matemático que posee magnitud (longitud) y dirección. Geométricamente, se representa como un segmento de recta orientado. Dos vectores son equivalentes si tienen la misma magnitud y dirección, independientemente de su punto de origen; por ello, los vectores también se denominan \textbf{vectores libres}.

🎧 Podcast: index\_definicion-de-vector

Tu navegador no soporta el elemento de audio.

\subsection{\texorpdfstring{Vector en la recta (\(\mathbb{R}\))}{Vector en la recta (\textbackslash mathbb\{R\})}}\label{vector-en-la-recta-mathbbr}

En la recta numérica, un vector se identifica con un único número real que indica su magnitud y sentido:

\[
\vec{v} = v_1 \in \mathbb{R}
\]

\subsubsection{🧮 Visualizador interactivo: Vector en la recta real}\label{visualizador-interactivo-vector-en-la-recta-real}

A continuación puede ingresar la componente \(v_1\) de un vector en \(\mathbb{R}\) y observar su representación geométrica junto con su magnitud.

\subsection{\texorpdfstring{Vector en el plano (\(\mathbb{R}^2\))}{Vector en el plano (\textbackslash mathbb\{R\}\^{}2)}}\label{vector-en-el-plano-mathbbr2}

En el plano cartesiano, un vector se expresa mediante un par ordenado de números reales:

\[
\vec{v} = (v_1, v_2) = v_1 \mathbf{i} + v_2 \mathbf{j}
\]

donde \(\mathbf{i} = (1, 0)\) y \(\mathbf{j} = (0, 1)\) son los vectores de la \textbf{base canónica}.

La magnitud (o norma) del vector es:

\[
\|\vec{v}\| = \sqrt{v_1^2 + v_2^2}
\]

\subsubsection{🧮 Visualizador interactivo: Vectores en el Plano}\label{visualizador-interactivo-vectores-en-el-plano}

Ingrese las componentes de uno o dos vectores para visualizarlos en el plano cartesiano, calcular su magnitud, dirección y operaciones.

\subsection{\texorpdfstring{Vector en el espacio (\(\mathbb{R}^3\))}{Vector en el espacio (\textbackslash mathbb\{R\}\^{}3)}}\label{vector-en-el-espacio-mathbbr3}

En el espacio tridimensional, un vector se representa por una terna ordenada:

\[
\vec{v} = (v_1, v_2, v_3) = v_1 \mathbf{i} + v_2 \mathbf{j} + v_3 \mathbf{k}
\]

donde \(\mathbf{i} = (1, 0, 0)\), \(\mathbf{j} = (0, 1, 0)\) y \(\mathbf{k} = (0, 0, 1)\).

La magnitud es:

\[
\|\vec{v}\| = \sqrt{v_1^2 + v_2^2 + v_3^2}
\]

Un vector de magnitud 1 se denomina \textbf{vector unitario}. Cualquier vector no nulo \(\vec{v}\) puede normalizarse dividiéndolo por su magnitud:

\[
\hat{v} = \frac{\vec{v}}{\|\vec{v}\|}
\]

\subsubsection{🎥 Video complementario: Vectores en el espacio}\label{video-complementario-vectores-en-el-espacio}

El siguiente video ilustra de forma visual el concepto de vectores en el espacio tridimensional, su representación gráfica y ejemplos prácticos.

\section{Operaciones vectoriales}\label{operaciones-vectoriales}

\subsection{Suma y resta}\label{suma-y-resta}

Dados \(\vec{u} = (u_1, u_2, u_3)\) y \(\vec{v} = (v_1, v_2, v_3)\):

\[
\vec{u} + \vec{v} = (u_1 + v_1,\; u_2 + v_2,\; u_3 + v_3)
\]

\[
\vec{u} - \vec{v} = (u_1 - v_1,\; u_2 - v_2,\; u_3 - v_3)
\]

Geométricamente, la suma se obtiene mediante la \textbf{regla del paralelogramo}.

\subsection{Multiplicación por escalar}\label{multiplicaciuxf3n-por-escalar}

Dado un escalar \(\lambda \in \mathbb{R}\):

\[
\lambda \vec{v} = (\lambda v_1,\; \lambda v_2,\; \lambda v_3)
\]

Si \(\lambda > 0\), el vector se alarga o acorta manteniendo la dirección. Si \(\lambda < 0\), invierte el sentido.

\section{Producto punto (producto escalar)}\label{producto-punto-producto-escalar}

\subsection{Definición algebraica}\label{definiciuxf3n-algebraica}

\[
\vec{u} \cdot \vec{v} = u_1 v_1 + u_2 v_2 + u_3 v_3
\]

🎧 Podcast: index\_definicion-algebraica

Tu navegador no soporta el elemento de audio.

\subsection{Definición geométrica}\label{definiciuxf3n-geomuxe9trica}

\[
\vec{u} \cdot \vec{v} = \|\vec{u}\| \|\vec{v}\| \cos\theta
\]

donde \(\theta\) es el ángulo entre los vectores (\(0 \leq \theta \leq \pi\)).

🎧 Podcast: index\_definicion-geometrica

Tu navegador no soporta el elemento de audio.

\subsection{Propiedades}\label{propiedades}

\begin{enumerate}
\def\labelenumi{\arabic{enumi}.}
\tightlist
\item
  \textbf{Conmutatividad}: \(\vec{u} \cdot \vec{v} = \vec{v} \cdot \vec{u}\)
\item
  \textbf{Distributividad}: \(\vec{u} \cdot (\vec{v} + \vec{w}) = \vec{u} \cdot \vec{v} + \vec{u} \cdot \vec{w}\)
\item
  \textbf{Asociatividad con escalares}: \((\lambda \vec{u}) \cdot \vec{v} = \lambda (\vec{u} \cdot \vec{v})\)
\item
  \textbf{Positividad}: \(\vec{v} \cdot \vec{v} = \|\vec{v}\|^2 \geq 0\), con igualdad solo si \(\vec{v} = \vec{0}\)
\end{enumerate}

🎧 Podcast: index\_propiedades

Tu navegador no soporta el elemento de audio.

\subsection{Aplicaciones del producto punto}\label{aplicaciones-del-producto-punto}

\textbf{Ángulo entre vectores}:

\[
\cos\theta = \frac{\vec{u} \cdot \vec{v}}{\|\vec{u}\| \|\vec{v}\|}
\]

\textbf{Vectores perpendiculares (ortogonales)}:

\[
\vec{u} \perp \vec{v} \iff \vec{u} \cdot \vec{v} = 0
\]

\textbf{Proyección escalar de \(\vec{u}\) sobre \(\vec{v}\)}:

\[
\text{comp}_{\vec{v}} \vec{u} = \frac{\vec{u} \cdot \vec{v}}{\|\vec{v}\|}
\]

\textbf{Proyección vectorial de \(\vec{u}\) sobre \(\vec{v}\)}:

\[
\text{proj}_{\vec{v}} \vec{u} = \left( \frac{\vec{u} \cdot \vec{v}}{\vec{v} \cdot \vec{v}} \right) \vec{v}
\]

\subsection{Ejemplo resuelto}\label{ejemplo-resuelto}

\begin{quote}
\textbf{Ejemplo 1.} Sean \(\vec{u} = (2, -1, 3)\) y \(\vec{v} = (1, 4, -2)\). Calcule \(\vec{u} \cdot \vec{v}\), el ángulo entre ellos, y la proyección de \(\vec{u}\) sobre \(\vec{v}\).
\end{quote}

\textbf{Solución.} El producto punto es:

\[
\vec{u} \cdot \vec{v} = 2(1) + (-1)(4) + 3(-2) = 2 - 4 - 6 = -8
\]

Las magnitudes son:

\[
\|\vec{u}\| = \sqrt{4 + 1 + 9} = \sqrt{14}, \qquad \|\vec{v}\| = \sqrt{1 + 16 + 4} = \sqrt{21}
\]

El ángulo satisface:

\[
\cos\theta = \frac{-8}{\sqrt{14}\sqrt{21}} = \frac{-8}{\sqrt{294}} \approx -0.466 \implies \theta \approx 117.8°
\]

La proyección vectorial es:

\[
\text{proj}_{\vec{v}} \vec{u} = \frac{-8}{21}(1, 4, -2) = \left(-\frac{8}{21}, -\frac{32}{21}, \frac{16}{21}\right)
\]

\section{Producto cruz (producto vectorial)}\label{producto-cruz-producto-vectorial}

El producto cruz está definido únicamente en \(\mathbb{R}^3\).

\subsection{Definición}\label{definiciuxf3n}

\[
\vec{u} \times \vec{v} = \begin{vmatrix}
\mathbf{i} & \mathbf{j} & \mathbf{k} \\
u_1 & u_2 & u_3 \\
v_1 & v_2 & v_3
\end{vmatrix} = (u_2 v_3 - u_3 v_2,\; u_3 v_1 - u_1 v_3,\; u_1 v_2 - u_2 v_1)
\]

🎧 Podcast: index\_definicion

Tu navegador no soporta el elemento de audio.

\subsection{Propiedades geométricas}\label{propiedades-geomuxe9tricas}

\begin{itemize}
\tightlist
\item
  El vector \(\vec{u} \times \vec{v}\) es perpendicular tanto a \(\vec{u}\) como a \(\vec{v}\).
\item
  Su magnitud es: \(\|\vec{u} \times \vec{v}\| = \|\vec{u}\| \|\vec{v}\| \sin\theta\), que corresponde al área del paralelogramo formado por \(\vec{u}\) y \(\vec{v}\).
\item
  La dirección sigue la \textbf{regla de la mano derecha}.
\end{itemize}

🎧 Podcast: index\_propiedades-geometricas

Tu navegador no soporta el elemento de audio.

\subsection{Propiedades algebraicas}\label{propiedades-algebraicas}

\begin{enumerate}
\def\labelenumi{\arabic{enumi}.}
\tightlist
\item
  \textbf{Anticonmutatividad}: \(\vec{u} \times \vec{v} = -(\vec{v} \times \vec{u})\)
\item
  \textbf{Distributividad}: \(\vec{u} \times (\vec{v} + \vec{w}) = \vec{u} \times \vec{v} + \vec{u} \times \vec{w}\)
\item
  \textbf{Asociatividad con escalares}: \((\lambda \vec{u}) \times \vec{v} = \lambda (\vec{u} \times \vec{v})\)
\item
  \textbf{Vector nulo}: \(\vec{u} \times \vec{u} = \vec{0}\)
\end{enumerate}

🎧 Podcast: index\_propiedades-algebraicas

Tu navegador no soporta el elemento de audio.

\subsection{Ejemplo resuelto}\label{ejemplo-resuelto-1}

\begin{quote}
\textbf{Ejemplo 2.} Calcule \(\vec{u} \times \vec{v}\) para \(\vec{u} = (1, 2, 3)\) y \(\vec{v} = (4, 5, 6)\). Verifique que es perpendicular a ambos.
\end{quote}

\textbf{Solución.}

\[
\vec{u} \times \vec{v} = \begin{vmatrix}
\mathbf{i} & \mathbf{j} & \mathbf{k} \\
1 & 2 & 3 \\
4 & 5 & 6
\end{vmatrix} = \mathbf{i}(12 - 15) - \mathbf{j}(6 - 12) + \mathbf{k}(5 - 8) = (-3, 6, -3)
\]

Verificación de perpendicularidad:

\[
\vec{u} \cdot (-3, 6, -3) = -3 + 12 - 9 = 0
\]

\[
\vec{v} \cdot (-3, 6, -3) = -12 + 30 - 18 = 0
\]

El área del paralelogramo formado por \(\vec{u}\) y \(\vec{v}\) es \(\|(-3, 6, -3)\| = \sqrt{9 + 36 + 9} = \sqrt{54} = 3\sqrt{6}\).

\section{Triple producto escalar}\label{triple-producto-escalar}

Dados tres vectores \(\vec{u}\), \(\vec{v}\), \(\vec{w}\) en \(\mathbb{R}^3\), el \textbf{triple producto escalar} es:

\[
\vec{u} \cdot (\vec{v} \times \vec{w}) = \begin{vmatrix}
u_1 & u_2 & u_3 \\
v_1 & v_2 & v_3 \\
w_1 & w_2 & w_3
\end{vmatrix}
\]

\subsection{Interpretación geométrica}\label{interpretaciuxf3n-geomuxe9trica}

El valor absoluto del triple producto escalar es el \textbf{volumen del paralelepípedo} formado por los tres vectores. Si el resultado es cero, los tres vectores son coplanares (linealmente dependientes).

\subsection{Ejemplo}\label{ejemplo}

\begin{quote}
\textbf{Ejemplo 3.} Calcule el volumen del paralelepípedo formado por \(\vec{u} = (1, 0, 0)\), \(\vec{v} = (0, 2, 0)\) y \(\vec{w} = (0, 0, 3)\).
\end{quote}

\textbf{Solución.}

\[
\vec{u} \cdot (\vec{v} \times \vec{w}) = \begin{vmatrix}
1 & 0 & 0 \\
0 & 2 & 0 \\
0 & 0 & 3
\end{vmatrix} = 1 \cdot 2 \cdot 3 = 6
\]

El volumen es \(|6| = 6\) unidades cúbicas (consistente con un prisma rectangular de dimensiones \(1 \times 2 \times 3\)).

\section{Visualización interactiva: vectores en 3D}\label{visualizaciuxf3n-interactiva-vectores-en-3d}

La siguiente aplicación interactiva le permite definir tres vectores en el espacio, visualizarlos en un entorno 3D orbitable y calcular todas las operaciones vectoriales fundamentales en tiempo real.

\section{Ejercicios propuestos}\label{ejercicios-propuestos}

\begin{enumerate}
\def\labelenumi{\arabic{enumi}.}
\tightlist
\item
  Dados \(\vec{a} = (3, -1, 2)\) y \(\vec{b} = (1, 2, -1)\), calcule \(\vec{a} + \vec{b}\), \(\vec{a} - \vec{b}\), \(2\vec{a} - 3\vec{b}\), \(\vec{a} \cdot \vec{b}\) y \(\vec{a} \times \vec{b}\).
\item
  Halle el ángulo entre los vectores \(\vec{u} = (1, 1, 1)\) y \(\vec{v} = (1, -1, 0)\).
\item
  Determine el valor de \(k\) para que los vectores \((2, k, 3)\) y \((1, 2, -1)\) sean perpendiculares.
\item
  Calcule el área del triángulo con vértices \(A = (1, 0, 0)\), \(B = (0, 2, 0)\) y \(C = (0, 0, 3)\).
\item
  Verifique que \((\vec{u} \times \vec{v}) \times \vec{w} \neq \vec{u} \times (\vec{v} \times \vec{w})\) en general, usando \(\vec{u} = (1, 0, 0)\), \(\vec{v} = (0, 1, 0)\), \(\vec{w} = (0, 0, 1)\).
\end{enumerate}

\chapter{Rectas y Planos en el Espacio}\label{rectas-planos}

🎧 Podcast del capítulo

Tu navegador no soporta el elemento de audio.

Los vectores son la herramienta fundamental para describir objetos geométricos en el espacio. En este capítulo se utilizan los conceptos vectoriales para estudiar rectas y planos, incluyendo sus ecuaciones, intersecciones, distancias y ángulos.

\section{Rectas en el espacio}\label{rectas-en-el-espacio}

Una recta en \(\mathbb{R}^3\) queda determinada por un punto \(P_0\) y un vector director \(\vec{v}\).

\subsection{Ecuación vectorial}\label{ecuaciuxf3n-vectorial}

Sea \(P_0 = (x_0, y_0, z_0)\) un punto de la recta y \(\vec{v} = (a, b, c)\) el vector director. Entonces cualquier punto \(P = (x, y, z)\) de la recta satisface:

\[
\vec{r} = \vec{r}_0 + t\vec{v}, \qquad t \in \mathbb{R}
\]

En coordenadas:

\[
(x, y, z) = (x_0, y_0, z_0) + t(a, b, c)
\]

\subsection{Ecuaciones paramétricas}\label{ecuaciones-paramuxe9tricas}

Descomponiendo la ecuación vectorial:

\[
\begin{aligned}
x &= x_0 + at \\
y &= y_0 + bt \\
z &= z_0 + ct
\end{aligned}, \qquad t \in \mathbb{R}
\]

\subsection{Ecuaciones simétricas}\label{ecuaciones-simuxe9tricas}

Si \(a, b, c \neq 0\), podemos despejar \(t\) en cada ecuación e igualar:

\[
\frac{x - x_0}{a} = \frac{y - y_0}{b} = \frac{z - z_0}{c}
\]

Si alguna componente del vector director es cero, la ecuación simétrica se adapta. Por ejemplo, si \(a = 0\):

\[
x = x_0, \qquad \frac{y - y_0}{b} = \frac{z - z_0}{c}
\]

\subsection{Recta que pasa por dos puntos}\label{recta-que-pasa-por-dos-puntos}

Dados \(P_1 = (x_1, y_1, z_1)\) y \(P_2 = (x_2, y_2, z_2)\), el vector director es \(\vec{v} = \overrightarrow{P_1 P_2} = (x_2 - x_1, y_2 - y_1, z_2 - z_1)\).

\subsection{Ejemplo resuelto}\label{ejemplo-resuelto-2}

\begin{quote}
\textbf{Ejemplo 1.} Halle las ecuaciones vectorial, paramétrica y simétrica de la recta que pasa por \(P_0 = (1, -2, 3)\) con dirección \(\vec{v} = (2, 1, -1)\).
\end{quote}

\textbf{Solución.}

\textbf{Vectorial:} \((x, y, z) = (1, -2, 3) + t(2, 1, -1)\)

\textbf{Paramétricas:}

\[
\begin{aligned}
x &= 1 + 2t \\
y &= -2 + t \\
z &= 3 - t
\end{aligned}
\]

\textbf{Simétricas:}

\[
\frac{x - 1}{2} = \frac{y + 2}{1} = \frac{z - 3}{-1}
\]

\subsubsection{🧮 Visualizador interactivo: Rectas en el espacio}\label{visualizador-interactivo-rectas-en-el-espacio}

Seleccione el método de definición de la recta, ingrese los datos y observe su representación geométrica junto con sus ecuaciones.

\section{Planos en el espacio}\label{planos-en-el-espacio}

Un plano queda determinado por un punto \(P_0\) y un vector normal \(\vec{n}\) perpendicular al plano.

\subsection{Ecuación punto-normal}\label{ecuaciuxf3n-punto-normal}

Si \(\vec{n} = (A, B, C)\) es el vector normal y \(P_0 = (x_0, y_0, z_0)\) es un punto del plano, entonces para cualquier punto \(P = (x, y, z)\) del plano, el vector \(\overrightarrow{P_0 P}\) es perpendicular a \(\vec{n}\):

\[
\vec{n} \cdot \overrightarrow{P_0 P} = 0
\]

Esto da:

\[
A(x - x_0) + B(y - y_0) + C(z - z_0) = 0
\]

\subsection{Ecuación general}\label{ecuaciuxf3n-general}

Expandiendo la ecuación punto-normal:

\[
Ax + By + Cz + D = 0
\]

donde \(D = -(Ax_0 + By_0 + Cz_0)\).

El vector normal es \(\vec{n} = (A, B, C)\). Si \((A, B, C) \neq (0, 0, 0)\), la ecuación representa un plano.

\subsection{Planos coordenados}\label{planos-coordenados}

\begin{longtable}[]{@{}lll@{}}
\toprule\noalign{}
Plano & Ecuación & Vector normal \\
\midrule\noalign{}
\endhead
\bottomrule\noalign{}
\endlastfoot
\(xy\) & \(z = 0\) & \((0, 0, 1)\) \\
\(xz\) & \(y = 0\) & \((0, 1, 0)\) \\
\(yz\) & \(x = 0\) & \((1, 0, 0)\) \\
\end{longtable}

\subsection{Interceptos}\label{interceptos}

Los \textbf{interceptos} del plano \(Ax + By + Cz = D\) (con \(D \neq 0\)) son:

\begin{itemize}
\tightlist
\item
  Intercepto \(x\): \(\left(\frac{D}{A}, 0, 0\right)\) (si \(A \neq 0\))
\item
  Intercepto \(y\): \(\left(0, \frac{D}{B}, 0\right)\) (si \(B \neq 0\))
\item
  Intercepto \(z\): \(\left(0, 0, \frac{D}{C}\right)\) (si \(C \neq 0\))
\end{itemize}

\subsection{Plano que pasa por tres puntos}\label{plano-que-pasa-por-tres-puntos}

Dados tres puntos no colineales \(P_1\), \(P_2\), \(P_3\), el vector normal se obtiene mediante el producto cruz:

\[
\vec{n} = \overrightarrow{P_1 P_2} \times \overrightarrow{P_1 P_3}
\]

\subsection{Ejemplo resuelto}\label{ejemplo-resuelto-3}

\begin{quote}
\textbf{Ejemplo 2.} Halle la ecuación del plano que pasa por \(P_1 = (1, 0, 0)\), \(P_2 = (0, 2, 0)\) y \(P_3 = (0, 0, 3)\).
\end{quote}

\textbf{Solución.} Los vectores en el plano son:

\[
\overrightarrow{P_1 P_2} = (-1, 2, 0), \qquad \overrightarrow{P_1 P_3} = (-1, 0, 3)
\]

El vector normal es:

\[
\vec{n} = (-1, 2, 0) \times (-1, 0, 3) = (6, 3, 2)
\]

Usando \(P_1\) en la ecuación punto-normal:

\[
6(x - 1) + 3(y - 0) + 2(z - 0) = 0 \implies 6x + 3y + 2z = 6
\]

O dividiendo entre 6:

\[
\frac{x}{1} + \frac{y}{2} + \frac{z}{3} = 1
\]

Esta forma muestra directamente los interceptos.

\subsubsection{🧮 Visualizador interactivo: Planos en el espacio}\label{visualizador-interactivo-planos-en-el-espacio}

Seleccione el método de definición del plano, ingrese los datos y observe su representación geométrica junto con sus ecuaciones.

\section{Distancias en el espacio}\label{distancias-en-el-espacio}

\subsection{Distancia de un punto a un plano}\label{distancia-de-un-punto-a-un-plano}

La distancia del punto \(Q = (x_1, y_1, z_1)\) al plano \(Ax + By + Cz + D = 0\) es:

\[
d(Q, \Pi) = \frac{|Ax_1 + By_1 + Cz_1 + D|}{\sqrt{A^2 + B^2 + C^2}}
\]

Esta fórmula se obtiene proyectando el vector \(\overrightarrow{P_0 Q}\) (donde \(P_0\) es cualquier punto del plano) sobre la dirección normal.

\subsection{Distancia de un punto a una recta}\label{distancia-de-un-punto-a-una-recta}

Dado un punto \(Q\) y una recta que pasa por \(P_0\) con dirección \(\vec{v}\), la distancia es:

\[
d(Q, L) = \frac{\|\overrightarrow{P_0 Q} \times \vec{v}\|}{\|\vec{v}\|}
\]

\subsection{Distancia entre dos rectas en el espacio}\label{distancia-entre-dos-rectas-en-el-espacio}

\begin{itemize}
\tightlist
\item
  Si las rectas \textbf{se intersectan}, la distancia es 0.
\item
  Si las rectas son \textbf{paralelas}, la distancia de cualquier punto de una recta a la otra es constante.
\item
  Si las rectas son \textbf{alabeadas} (no paralelas y no se intersectan), la distancia es:
\end{itemize}

\[
d(L_1, L_2) = \frac{|(\vec{r}_2 - \vec{r}_1) \cdot (\vec{v}_1 \times \vec{v}_2)|}{\|\vec{v}_1 \times \vec{v}_2\|}
\]

donde \(\vec{r}_1, \vec{v}_1\) y \(\vec{r}_2, \vec{v}_2\) son puntos y vectores directores de cada recta.

\subsection{Ejemplo resuelto}\label{ejemplo-resuelto-4}

\begin{quote}
\textbf{Ejemplo 3.} Calcule la distancia del punto \(Q = (1, 2, 3)\) al plano \(2x - y + 2z - 4 = 0\).
\end{quote}

\textbf{Solución.}

\[
d = \frac{|2(1) - 1(2) + 2(3) - 4|}{\sqrt{4 + 1 + 4}} = \frac{|2 - 2 + 6 - 4|}{3} = \frac{2}{3}
\]

\subsubsection{🧮 Calculadora interactiva: Distancias en el espacio}\label{calculadora-interactiva-distancias-en-el-espacio}

Seleccione el tipo de distancia o magnitud que desea calcular. Ingrese los datos correspondientes y observe el resultado numérico junto con su representación geométrica.

\section{Ángulos en el espacio}\label{uxe1ngulos-en-el-espacio}

\subsection{Ángulo entre dos rectas}\label{uxe1ngulo-entre-dos-rectas}

Si las rectas tienen vectores directores \(\vec{v}_1\) y \(\vec{v}_2\):

\[
\cos\theta = \frac{|\vec{v}_1 \cdot \vec{v}_2|}{\|\vec{v}_1\| \|\vec{v}_2\|}
\]

Se toma el valor absoluto para obtener el ángulo agudo.

\subsection{Ángulo entre dos planos}\label{uxe1ngulo-entre-dos-planos}

Si los planos tienen vectores normales \(\vec{n}_1\) y \(\vec{n}_2\):

\[
\cos\theta = \frac{|\vec{n}_1 \cdot \vec{n}_2|}{\|\vec{n}_1\| \|\vec{n}_2\|}
\]

\subsection{Ángulo entre una recta y un plano}\label{uxe1ngulo-entre-una-recta-y-un-plano}

Si la recta tiene vector director \(\vec{v}\) y el plano tiene vector normal \(\vec{n}\):

\[
\sin\theta = \frac{|\vec{v} \cdot \vec{n}|}{\|\vec{v}\| \|\vec{n}\|}
\]

donde \(\theta\) es el ángulo entre la recta y su proyección sobre el plano.

\subsubsection{🧮 Calculadora interactiva: Ángulos en el espacio}\label{calculadora-interactiva-uxe1ngulos-en-el-espacio}

Calcule el ángulo entre vectores, rectas, planos, y entre rectas y planos. Observe la representación geométrica y el resultado en grados y radianes.

\section{Intersecciones}\label{intersecciones}

\subsection{Intersección de una recta y un plano}\label{intersecciuxf3n-de-una-recta-y-un-plano}

Sustituya las ecuaciones paramétricas de la recta en la ecuación del plano y resuelva para \(t\).

\begin{itemize}
\tightlist
\item
  Si obtiene un único valor de \(t\): la recta corta al plano en un punto.
\item
  Si la ecuación es una identidad (\(0 = 0\)): la recta está contenida en el plano.
\item
  Si la ecuación es imposible (\(0 = k\) con \(k \neq 0\)): la recta es paralela al plano (no se intersectan).
\end{itemize}

\subsection{Intersección de dos planos}\label{intersecciuxf3n-de-dos-planos}

Dos planos no paralelos se intersectan en una \textbf{recta}. La dirección de esta recta es perpendicular a ambos vectores normales:

\[
\vec{v} = \vec{n}_1 \times \vec{n}_2
\]

Para hallar un punto de la recta de intersección, se resuelve el sistema de las dos ecuaciones de los planos (imponiendo una condición adicional, por ejemplo \(z = 0\) si es posible).

\subsection{Intersección de tres planos}\label{intersecciuxf3n-de-tres-planos}

Tres planos pueden:

\begin{enumerate}
\def\labelenumi{\arabic{enumi}.}
\tightlist
\item
  \textbf{Intersectarse en un único punto}: el sistema tiene solución única.
\item
  \textbf{Intersectarse en una recta común}: los tres planos pertenecen a un haz.
\item
  \textbf{No tener intersección común}: el sistema es incompatible.
\item
  \textbf{Coincidir dos a dos}: dos o más planos son el mismo.
\end{enumerate}

\subsubsection{🧮 Visualizador interactivo: Intersecciones en el espacio}\label{visualizador-interactivo-intersecciones-en-el-espacio}

Calcule y visualice la intersección entre rectas, planos, y combinaciones de ambos. La aplicación resuelve el sistema correspondiente y muestra el punto, la recta o la configuración resultante.

\section{Ejercicios propuestos}\label{ejercicios-propuestos-1}

\begin{enumerate}
\def\labelenumi{\arabic{enumi}.}
\tightlist
\item
  Halle las ecuaciones paramétricas y simétricas de la recta que pasa por \((2, -1, 4)\) y es paralela al vector \((3, 2, -2)\).
\item
  Determine la ecuación del plano que pasa por \((1, 1, 1)\) y es perpendicular a la recta \(x = 2t, y = 1 - t, z = 3 + 2t\).
\item
  Calcule la distancia entre las rectas alabeadas \(L_1: x = 1 + t, y = 2 - t, z = 3t\) y \(L_2: x = 2s, y = 1 + s, z = -1 + 2s\).
\item
  Halle el ángulo entre los planos \(x + y + z = 1\) y \(x - y + z = 2\).
\item
  Determine si la recta \(\frac{x-1}{2} = \frac{y+1}{-1} = \frac{z-3}{1}\) intersecta al plano \(x + 2y - z = 5\). En caso afirmativo, halle el punto de intersección.
\end{enumerate}

\chapter{Sistemas de Coordenadas y Parametrización}\label{coordenadas}

🎧 Podcast del capítulo

Tu navegador no soporta el elemento de audio.

Además de las coordenadas cartesianas, existen otros sistemas de coordenadas que simplifican enormemente la descripción de ciertas curvas y superficies. En este capítulo se estudian los sistemas polares, cilíndricos y esféricos, así como la parametrización de curvas en dos y tres dimensiones.

\section{Coordenadas polares en el plano}\label{coordenadas-polares-en-el-plano}

En el sistema de \textbf{coordenadas polares}, un punto \(P\) del plano se describe mediante la distancia \(r\) al origen y el ángulo \(\theta\) que forma el segmento \(OP\) con el eje \(x\) positivo.

\subsection{Relación con coordenadas cartesianas}\label{relaciuxf3n-con-coordenadas-cartesianas}

\[
\begin{aligned}
x &= r \cos\theta \\
y &= r \sin\theta
\end{aligned}
\]

Inversamente:

\[
\begin{aligned}
r &= \sqrt{x^2 + y^2} \\
\theta &= \arctan\left(\frac{y}{x}\right)
\end{aligned}
\]

\subsection{Curvas notables en coordenadas polares}\label{curvas-notables-en-coordenadas-polares}

\begin{longtable}[]{@{}
  >{\raggedright\arraybackslash}p{(\columnwidth - 4\tabcolsep) * \real{0.1944}}
  >{\raggedright\arraybackslash}p{(\columnwidth - 4\tabcolsep) * \real{0.4444}}
  >{\raggedright\arraybackslash}p{(\columnwidth - 4\tabcolsep) * \real{0.3611}}@{}}
\toprule\noalign{}
\begin{minipage}[b]{\linewidth}\raggedright
Curva
\end{minipage} & \begin{minipage}[b]{\linewidth}\raggedright
Ecuación polar
\end{minipage} & \begin{minipage}[b]{\linewidth}\raggedright
Descripción
\end{minipage} \\
\midrule\noalign{}
\endhead
\bottomrule\noalign{}
\endlastfoot
Circunferencia & \(r = a\) & Centro en el origen, radio \(a\) \\
Circunferencia & \(r = 2a\cos\theta\) & Centro en \((a, 0)\), radio \(a\) \\
Recta por el origen & \(\theta = \alpha\) & Ángulo constante \\
Cardioide & \(r = a(1 + \cos\theta)\) & Corazón con cúspide en el origen \\
Rosa de 4 pétalos & \(r = a\cos(2\theta)\) & Pétalos simétricos \\
Rosa de 3 pétalos & \(r = a\cos(3\theta)\) & Pétalos alternados \\
Lemniscata & \(r^2 = a^2\cos(2\theta)\) & Figura en forma de 8 \\
Espiral de Arquímedes & \(r = a\theta\) & Espiral que se abre uniformemente \\
\end{longtable}

\subsection{Ejemplo resuelto}\label{ejemplo-resuelto-5}

\begin{quote}
\textbf{Ejemplo 1.} Convierta la ecuación cartesiana \(x^2 + y^2 = 2x\) a coordenadas polares y grafique.
\end{quote}

\textbf{Solución.} Sustituyendo \(x = r\cos\theta\), \(y = r\sin\theta\):

\[
r^2 = 2r\cos\theta \implies r = 2\cos\theta
\]

Es una circunferencia de radio 1 centrada en \((1, 0)\).

\begin{figure}

{\centering \includegraphics[width=1\linewidth]{Superficies3d_files/figure-latex/cardioide-1} 

}

\caption{Cardioide $r = 1 + \cos\theta$ en coordenadas polares.}\label{fig:cardioide}
\end{figure}

\section{Coordenadas cilíndricas}\label{coordenadas-ciluxedndricas}

El sistema de \textbf{coordenadas cilíndricas} extiende las coordenadas polares al espacio tridimensional agregando la coordenada vertical \(z\).

\subsection{Definición}\label{definiciuxf3n-1}

Un punto \(P\) en el espacio se describe por \((r, \theta, z)\) donde:

\begin{itemize}
\tightlist
\item
  \(r\) es la distancia del punto al eje \(z\).
\item
  \(\theta\) es el ángulo alrededor del eje \(z\).
\item
  \(z\) es la altura sobre el plano \(xy\).
\end{itemize}

🎧 Podcast: 02\_coordenadas\_parametrizacion\_definicion

Tu navegador no soporta el elemento de audio.

\subsection{Conversión a cartesianas}\label{conversiuxf3n-a-cartesianas}

\[
\begin{aligned}
x &= r \cos\theta \\
y &= r \sin\theta \\
z &= z
\end{aligned}
\]

\subsection{Superficies de coordenadas constantes}\label{superficies-de-coordenadas-constantes}

\begin{longtable}[]{@{}lll@{}}
\toprule\noalign{}
Superficie & Ecuación & Descripción \\
\midrule\noalign{}
\endhead
\bottomrule\noalign{}
\endlastfoot
Cilindro circular & \(r = a\) & Eje \(z\), radio \(a\) \\
Plano vertical & \(\theta = \alpha\) & Contiene el eje \(z\) \\
Plano horizontal & \(z = k\) & Paralelo al plano \(xy\) \\
\end{longtable}

\subsection{Jacobiano}\label{jacobiano}

El elemento de volumen en coordenadas cilíndricas es:

\[
dV = r \, dr \, d\theta \, dz
\]

El factor \(r\) es el \textbf{jacobiano} de la transformación.

\subsubsection{🧮 Visualizador interactivo: coordenadas cilíndricas en 3D}\label{visualizador-interactivo-coordenadas-ciluxedndricas-en-3d}

Explore superficies definidas por \(r = f(\theta, z)\) o trace curvas paramétricas expresadas en coordenadas cilíndricas \((r(t), \theta(t), z(t))\). La aplicación calcula la longitud de arco para curvas.

\section{Coordenadas esféricas}\label{coordenadas-esfuxe9ricas}

En el sistema de \textbf{coordenadas esféricas}, un punto se describe por \((\rho, \phi, \theta)\) donde:

\begin{itemize}
\tightlist
\item
  \(\rho\) es la distancia al origen (\(\rho \geq 0\)).
\item
  \(\phi\) es el ángulo desde el eje \(z\) positivo (\(0 \leq \phi \leq \pi\)).
\item
  \(\theta\) es el ángulo azimutal en el plano \(xy\) desde el eje \(x\) (\(0 \leq \theta < 2\pi\)).
\end{itemize}

\subsection{Conversión a cartesianas}\label{conversiuxf3n-a-cartesianas-1}

\[
\begin{aligned}
x &= \rho \sin\phi \cos\theta \\
y &= \rho \sin\phi \sin\theta \\
z &= \rho \cos\phi
\end{aligned}
\]

\subsection{Conversión a cilíndricas}\label{conversiuxf3n-a-ciluxedndricas}

\[
r = \rho \sin\phi, \qquad z = \rho \cos\phi, \qquad \theta = \theta
\]

\subsection{Superficies de coordenadas constantes}\label{superficies-de-coordenadas-constantes-1}

\begin{longtable}[]{@{}lll@{}}
\toprule\noalign{}
Superficie & Ecuación & Descripción \\
\midrule\noalign{}
\endhead
\bottomrule\noalign{}
\endlastfoot
Esfera & \(\rho = a\) & Centro en el origen, radio \(a\) \\
Cono & \(\phi = \alpha\) & Vértice en el origen, abertura \(2\alpha\) \\
Plano vertical & \(\theta = \beta\) & Contiene el eje \(z\) \\
\end{longtable}

\subsection{Jacobiano}\label{jacobiano-1}

El elemento de volumen en coordenadas esféricas es:

\[
dV = \rho^2 \sin\phi \, d\rho \, d\phi \, d\theta
\]

\subsection{Ejemplo resuelto}\label{ejemplo-resuelto-6}

\begin{quote}
\textbf{Ejemplo 2.} El punto con coordenadas esféricas \((\rho, \phi, \theta) = (2, \pi/4, \pi/3)\) tiene coordenadas cartesianas:
\end{quote}

\textbf{Solución.}

\[
\begin{aligned}
x &= 2 \sin(\pi/4) \cos(\pi/3) = 2 \cdot \frac{\sqrt{2}}{2} \cdot \frac{1}{2} = \frac{\sqrt{2}}{2} \\
y &= 2 \sin(\pi/4) \sin(\pi/3) = 2 \cdot \frac{\sqrt{2}}{2} \cdot \frac{\sqrt{3}}{2} = \frac{\sqrt{6}}{2} \\
z &= 2 \cos(\pi/4) = 2 \cdot \frac{\sqrt{2}}{2} = \sqrt{2}
\end{aligned}
\]

\label{fig:esfera-coord}Esfera de radio 2 visualizada con coordenadas esféricas.

\subsubsection{🧮 Visualizador interactivo: coordenadas esféricas en 3D}\label{visualizador-interactivo-coordenadas-esfuxe9ricas-en-3d}

Explore superficies definidas por \(\rho = f(\phi, \theta)\) o trace curvas paramétricas expresadas en coordenadas esféricas \((\rho(t), \phi(t), \theta(t))\). La aplicación calcula la longitud de arco para curvas.

\subsubsection{🧮 Visualizador interactivo: intersección de tres superficies y cálculo de volumen}\label{visualizador-interactivo-intersecciuxf3n-de-tres-superficies-y-cuxe1lculo-de-volumen}

Explore la región sólida formada por la intersección de hasta tres superficies. Visualice las superficies, la nube de puntos del sólido resultante y calcule su volumen por el método de Monte Carlo.

\section{Parametrización de curvas}\label{parametrizaciuxf3n-de-curvas}

Una curva en el plano o en el espacio puede describirse mediante una función vectorial de un parámetro.

\subsection{Curvas paramétricas en el plano}\label{curvas-paramuxe9tricas-en-el-plano}

Una curva en \(\mathbb{R}^2\) se describe como:

\[
\vec{r}(t) = \bigl( x(t),\; y(t) \bigr), \qquad t \in [a, b]
\]

\subsubsection{Ejemplos clásicos}\label{ejemplos-cluxe1sicos}

\textbf{Circunferencia de radio \(R\):}

\[
x = R\cos t, \quad y = R\sin t, \quad t \in [0, 2\pi]
\]

\textbf{Elipse de semiejes \(a\) y \(b\):}

\[
x = a\cos t, \quad y = b\sin t, \quad t \in [0, 2\pi]
\]

\textbf{Cicloide} (trazada por un punto de una rueda de radio \(R\) que rueda sin deslizar):

\[
x = R(t - \sin t), \quad y = R(1 - \cos t), \quad t \in [0, 2\pi]
\]

\textbf{Astroide}:

\[
x = a\cos^3 t, \quad y = a\sin^3 t, \quad t \in [0, 2\pi]
\]

\textbf{Epicicloide} (trazada por un punto de una circunferencia de radio \(r\) que rueda alrededor de otra de radio \(R\)):

\[
\begin{aligned}
x &= (R + r)\cos t - r\cos\left(\frac{R + r}{r}t\right) \\
y &= (R + r)\sin t - r\sin\left(\frac{R + r}{r}t\right)
\end{aligned}
\]

\begin{figure}

{\centering \includegraphics[width=1\linewidth]{Superficies3d_files/figure-latex/cicloide-1} 

}

\caption{Cicloide generada por una rueda de radio $R = 1$.}\label{fig:cicloide}
\end{figure}

\subsection{Curvas paramétricas en el espacio}\label{curvas-paramuxe9tricas-en-el-espacio}

Una curva en \(\mathbb{R}^3\) se describe como:

\[
\vec{r}(t) = \bigl( x(t),\; y(t),\; z(t) \bigr), \quad t \in [a, b]
\]

\subsubsection{Ejemplos clásicos}\label{ejemplos-cluxe1sicos-1}

\textbf{Hélice circular} (resorte):

\[
x = a\cos t, \quad y = a\sin t, \quad z = ct
\]

\textbf{Curva de Viviani} (intersección de una esfera y un cilindro):

\[
x = a(1 + \cos t), \quad y = a\sin t, \quad z = 2a\sin(t/2)
\]

\label{fig:helice}Hélice circular de radio \(a = 2\) y paso \(c = 0.5\).

\section{Longitud de arco}\label{longitud-de-arco}

La \textbf{longitud de arco} de una curva paramétrica desde \(t = a\) hasta \(t = b\) es:

\[
L = \int_a^b \|\vec{r}'(t)\| \, dt = \int_a^b \sqrt{[x'(t)]^2 + [y'(t)]^2 + [z'(t)]^2} \, dt
\]

\subsection{Ejemplo resuelto}\label{ejemplo-resuelto-7}

\begin{quote}
\textbf{Ejemplo 3.} Calcule la longitud de la hélice \(\vec{r}(t) = (a\cos t, a\sin t, ct)\) para \(t \in [0, 2\pi]\).
\end{quote}

\textbf{Solución.} La derivada es:

\[
\vec{r}'(t) = (-a\sin t, a\cos t, c)
\]

Su magnitud es:

\[
\|\vec{r}'(t)\| = \sqrt{a^2\sin^2 t + a^2\cos^2 t + c^2} = \sqrt{a^2 + c^2}
\]

Por tanto:

\[
L = \int_0^{2\pi} \sqrt{a^2 + c^2} \, dt = 2\pi\sqrt{a^2 + c^2}
\]

\subsubsection{🧮 Visualizador interactivo: coordenadas polares}\label{visualizador-interactivo-coordenadas-polares}

Introduzca una ecuación \(r = f(\theta)\) y el rango angular para graficar la curva en coordenadas polares. La aplicación convierte automáticamente a cartesianas y dibuja la figura.

\subsubsection{🧮 Visualizador interactivo: curvas paramétricas cartesianas}\label{visualizador-interactivo-curvas-paramuxe9tricas-cartesianas}

Introduzca las funciones \(x(t)\) y \(y(t)\) y el intervalo del parámetro \(t\). La aplicación dibuja la curva paramétrica en el plano cartesiano.

\section{Ejercicios propuestos}\label{ejercicios-propuestos-2}

\begin{enumerate}
\def\labelenumi{\arabic{enumi}.}
\tightlist
\item
  Convierta el punto con coordenadas polares \((r, \theta) = (3, \pi/6)\) a coordenadas cartesianas.
\item
  Halle la ecuación polar de la recta \(y = x\).
\item
  Calcule la longitud de arco de la astroide \(x = \cos^3 t\), \(y = \sin^3 t\) para \(t \in [0, 2\pi]\).
\item
  Parametrice la elipse \(\frac{x^2}{4} + \frac{y^2}{9} = 1\) y calcule su longitud de arco (expresada como una integral elíptica).
\item
  Un punto tiene coordenadas cartesianas \((1, 1, \sqrt{2})\). Halle sus coordenadas cilíndricas y esféricas.
\item
  Describa la superficie dada en coordenadas esféricas por \(\rho = 2\cos\phi\) y conviértala a coordenadas cartesianas.
\end{enumerate}

\chapter{Funciones de Varias Variables y Derivadas Parciales}\label{derivadas-parciales}

🎧 Podcast del capítulo

Tu navegador no soporta el elemento de audio.

El cálculo diferencial de funciones de varias variables es la base del análisis de superficies y campos escalares. En este capítulo se desarrollan las herramientas fundamentales: límites, continuidad, derivadas parciales, plano tangente, regla de la cadena, derivada direccional y gradiente.

\section{Funciones de dos y tres variables}\label{funciones-de-dos-y-tres-variables}

\subsection{Definición}\label{definiciuxf3n-2}

Una \textbf{función de dos variables} es una regla que asigna a cada par ordenado \((x, y)\) en un conjunto \(D \subseteq \mathbb{R}^2\) un único número real \(z = f(x, y)\).

De manera análoga, una \textbf{función de tres variables} asigna a cada terna \((x, y, z) \in D \subseteq \mathbb{R}^3\) un valor \(w = f(x, y, z)\).

🎧 Podcast: 03\_derivadas\_parciales\_definicion

Tu navegador no soporta el elemento de audio.

\subsection{Dominio y rango}\label{dominio-y-rango}

El \textbf{dominio} de \(f\) es el conjunto de todos los puntos donde la función está definida. El \textbf{rango} es el conjunto de todos los valores posibles de \(f\).

\subsection{Ejemplos}\label{ejemplos}

\begin{longtable}[]{@{}
  >{\raggedright\arraybackslash}p{(\columnwidth - 4\tabcolsep) * \real{0.3600}}
  >{\raggedright\arraybackslash}p{(\columnwidth - 4\tabcolsep) * \real{0.3600}}
  >{\raggedright\arraybackslash}p{(\columnwidth - 4\tabcolsep) * \real{0.2800}}@{}}
\toprule\noalign{}
\begin{minipage}[b]{\linewidth}\raggedright
Función
\end{minipage} & \begin{minipage}[b]{\linewidth}\raggedright
Dominio
\end{minipage} & \begin{minipage}[b]{\linewidth}\raggedright
Rango
\end{minipage} \\
\midrule\noalign{}
\endhead
\bottomrule\noalign{}
\endlastfoot
\(f(x, y) = x^2 + y^2\) & \(\mathbb{R}^2\) & \([0, \infty)\) \\
\(f(x, y) = \sqrt{1 - x^2 - y^2}\) & \(x^2 + y^2 \leq 1\) & \([0, 1]\) \\
\(f(x, y) = \frac{1}{x - y}\) & \(x \neq y\) & \(\mathbb{R} \setminus \{0\}\) \\
\(f(x, y, z) = x^2 + y^2 + z^2\) & \(\mathbb{R}^3\) & \([0, \infty)\) \\
\end{longtable}

\section{Gráfica de una función de dos variables}\label{gruxe1fica-de-una-funciuxf3n-de-dos-variables}

La gráfica de \(z = f(x, y)\) es el conjunto de puntos \((x, y, z) \in \mathbb{R}^3\) tal que \(z = f(x, y)\). Geométricamente, es una \textbf{superficie} en el espacio tridimensional.

\subsection{Ejemplo}\label{ejemplo-1}

La gráfica de \(f(x, y) = x^2 + y^2\) es un \textbf{paraboloide elíptico} que abre hacia arriba con vértice en el origen.

\label{fig:paraboloide-ejemplo}Paraboloide elíptico \(z = x^2 + y^2\).

\section{Curvas de nivel}\label{curvas-de-nivel}

Las \textbf{curvas de nivel} de una función \(f(x, y)\) son las curvas en el plano \(xy\) definidas por \(f(x, y) = k\), donde \(k\) es una constante. Proporcionan una representación bidimensional de la función, análoga a los mapas topográficos.

\subsection{Ejemplos de curvas de nivel}\label{ejemplos-de-curvas-de-nivel}

\begin{longtable}[]{@{}lll@{}}
\toprule\noalign{}
Función & Curvas de nivel & Descripción \\
\midrule\noalign{}
\endhead
\bottomrule\noalign{}
\endlastfoot
\(f(x, y) = x^2 + y^2\) & \(x^2 + y^2 = k\) & Círculos concéntricos \\
\(f(x, y) = x^2 - y^2\) & \(x^2 - y^2 = k\) & Hipérbolas \\
\(f(x, y) = x + y\) & \(x + y = k\) & Rectas paralelas \\
\end{longtable}

\begin{figure}

{\centering \includegraphics[width=1\linewidth]{Superficies3d_files/figure-latex/curvas-nivel-1} 

}

\caption{Curvas de nivel del paraboloide $z = x^2 + y^2$.}\label{fig:curvas-nivel}
\end{figure}

\section{Límites y continuidad}\label{luxedmites-y-continuidad}

\subsection{Límite de una función de dos variables}\label{luxedmite-de-una-funciuxf3n-de-dos-variables}

Decimos que \(\lim_{(x,y) \to (a,b)} f(x, y) = L\) si para todo \(\varepsilon > 0\) existe \(\delta > 0\) tal que:

\[
0 < \sqrt{(x-a)^2 + (y-b)^2} < \delta \implies |f(x, y) - L| < \varepsilon
\]

La dificultad principal es que \((x, y)\) puede aproximarse a \((a, b)\) por \textbf{infinitos caminos}.

\subsection{Criterio para inexistencia de límite}\label{criterio-para-inexistencia-de-luxedmite}

Si al aproximarnos por dos caminos distintos obtenemos valores límite diferentes, entonces el límite \textbf{no existe}.

\subsection{Ejemplo resuelto}\label{ejemplo-resuelto-8}

\begin{quote}
\textbf{Ejemplo 1.} Demuestre que \(\lim_{(x,y) \to (0,0)} \frac{xy}{x^2 + y^2}\) no existe.
\end{quote}

\textbf{Solución.} Aproximando por el eje \(x\) (\(y = 0\)):

\[
\lim_{x \to 0} \frac{x \cdot 0}{x^2 + 0} = 0
\]

Aproximando por la recta \(y = x\):

\[
\lim_{x \to 0} \frac{x \cdot x}{x^2 + x^2} = \lim_{x \to 0} \frac{x^2}{2x^2} = \frac{1}{2}
\]

Como \(0 \neq \frac{1}{2}\), el límite no existe.

\subsection{Continuidad}\label{continuidad}

Una función \(f\) es \textbf{continua} en \((a, b)\) si:

\begin{enumerate}
\def\labelenumi{\arabic{enumi}.}
\tightlist
\item
  \(f(a, b)\) está definida.
\item
  \(\lim_{(x,y) \to (a,b)} f(x, y)\) existe.
\item
  \(\lim_{(x,y) \to (a,b)} f(x, y) = f(a, b)\).
\end{enumerate}

\section{Derivadas parciales}\label{derivadas-parciales-1}

\subsection{Definición}\label{definiciuxf3n-3}

La \textbf{derivada parcial} de \(f\) respecto a \(x\) en \((a, b)\) se define como:

\[
f_x(a, b) = \lim_{h \to 0} \frac{f(a + h, b) - f(a, b)}{h}
\]

De manera análoga:

\[
f_y(a, b) = \lim_{h \to 0} \frac{f(a, b + h) - f(a, b)}{h}
\]

En la práctica, calculamos \(f_x\) derivando \(f\) respecto a \(x\) tratando a \(y\) como constante, y viceversa.

🎧 Podcast: 03\_derivadas\_parciales\_definicion\_1

Tu navegador no soporta el elemento de audio.

\subsection{Notaciones}\label{notaciones}

\[
f_x = \frac{\partial f}{\partial x} = \partial_x f, \qquad
f_y = \frac{\partial f}{\partial y} = \partial_y f
\]

\subsection{Interpretación geométrica}\label{interpretaciuxf3n-geomuxe9trica-1}

\begin{itemize}
\tightlist
\item
  \(f_x(a, b)\) es la pendiente de la recta tangente a la curva de intersección de la superficie \(z = f(x, y)\) con el plano vertical \(y = b\).
\item
  \(f_y(a, b)\) es la pendiente de la recta tangente a la curva de intersección con el plano vertical \(x = a\).
\end{itemize}

\subsection{Ejemplo resuelto}\label{ejemplo-resuelto-9}

\begin{quote}
\textbf{Ejemplo 2.} Calcule las derivadas parciales de \(f(x, y) = x^3 y^2 + 2xy^4\).
\end{quote}

\textbf{Solución.}

\[
f_x = 3x^2 y^2 + 2y^4
\]

\[
f_y = 2x^3 y + 8xy^3
\]

\section{Herramienta de cálculo: Derivadas parciales paso a paso}\label{herramienta-de-cuxe1lculo-derivadas-parciales-paso-a-paso}

A continuación se presenta una herramienta de cálculo simbólico que permite obtener las derivadas parciales de cualquier función polinómica. El usuario puede modificar la función en el código para calcular sus propias derivadas.

\subsection{\texorpdfstring{Ejemplo: \(f(x, y) = x^3 \sin(y) + e^{xy}\)}{Ejemplo: f(x, y) = x\^{}3 \textbackslash sin(y) + e\^{}\{xy\}}}\label{ejemplo-fx-y-x3-siny-exy}

\begin{Shaded}
\begin{Highlighting}[]
\FunctionTok{library}\NormalTok{(Deriv)}

\CommentTok{\# Defina aquí su función}
\NormalTok{f }\OtherTok{\textless{}{-}} \ControlFlowTok{function}\NormalTok{(x, y) x}\SpecialCharTok{\^{}}\DecValTok{3} \SpecialCharTok{*} \FunctionTok{sin}\NormalTok{(y) }\SpecialCharTok{+} \FunctionTok{exp}\NormalTok{(x }\SpecialCharTok{*}\NormalTok{ y)}

\CommentTok{\# Cálculo simbólico}
\NormalTok{fx }\OtherTok{\textless{}{-}} \FunctionTok{Deriv}\NormalTok{(f, }\StringTok{"x"}\NormalTok{)}
\NormalTok{fy }\OtherTok{\textless{}{-}} \FunctionTok{Deriv}\NormalTok{(f, }\StringTok{"y"}\NormalTok{)}
\NormalTok{fxx }\OtherTok{\textless{}{-}} \FunctionTok{Deriv}\NormalTok{(fx, }\StringTok{"x"}\NormalTok{)}
\NormalTok{fyy }\OtherTok{\textless{}{-}} \FunctionTok{Deriv}\NormalTok{(fy, }\StringTok{"y"}\NormalTok{)}
\NormalTok{fxy }\OtherTok{\textless{}{-}} \FunctionTok{Deriv}\NormalTok{(fx, }\StringTok{"y"}\NormalTok{)}

\FunctionTok{cat}\NormalTok{(}\StringTok{"Función original: f(x,y) = x³·sin(y) + e\^{}(x·y)}\SpecialCharTok{\textbackslash{}n\textbackslash{}n}\StringTok{"}\NormalTok{)}
\end{Highlighting}
\end{Shaded}

\begin{verbatim}
## Función original: f(x,y) = x³·sin(y) + e^(x·y)
\end{verbatim}

\begin{Shaded}
\begin{Highlighting}[]
\FunctionTok{cat}\NormalTok{(}\StringTok{"∂f/∂x ="}\NormalTok{, }\FunctionTok{as.character}\NormalTok{(}\FunctionTok{body}\NormalTok{(fx)), }\StringTok{"}\SpecialCharTok{\textbackslash{}n}\StringTok{"}\NormalTok{)}
\end{Highlighting}
\end{Shaded}

\begin{verbatim}
## ∂f/∂x = + 3 * (x^2 * sin(y)) y * exp(x * y)
\end{verbatim}

\begin{Shaded}
\begin{Highlighting}[]
\FunctionTok{cat}\NormalTok{(}\StringTok{"∂f/∂y ="}\NormalTok{, }\FunctionTok{as.character}\NormalTok{(}\FunctionTok{body}\NormalTok{(fy)), }\StringTok{"}\SpecialCharTok{\textbackslash{}n}\StringTok{"}\NormalTok{)}
\end{Highlighting}
\end{Shaded}

\begin{verbatim}
## ∂f/∂y = * x exp(x * y) + x^2 * cos(y)
\end{verbatim}

\begin{Shaded}
\begin{Highlighting}[]
\FunctionTok{cat}\NormalTok{(}\StringTok{"∂²f/∂x² ="}\NormalTok{, }\FunctionTok{as.character}\NormalTok{(}\FunctionTok{body}\NormalTok{(fxx)), }\StringTok{"}\SpecialCharTok{\textbackslash{}n}\StringTok{"}\NormalTok{)}
\end{Highlighting}
\end{Shaded}

\begin{verbatim}
## ∂²f/∂x² = + 6 * (x * sin(y)) y^2 * exp(x * y)
\end{verbatim}

\begin{Shaded}
\begin{Highlighting}[]
\FunctionTok{cat}\NormalTok{(}\StringTok{"∂²f/∂y² ="}\NormalTok{, }\FunctionTok{as.character}\NormalTok{(}\FunctionTok{body}\NormalTok{(fyy)), }\StringTok{"}\SpecialCharTok{\textbackslash{}n}\StringTok{"}\NormalTok{)}
\end{Highlighting}
\end{Shaded}

\begin{verbatim}
## ∂²f/∂y² = * x^2 exp(x * y) - x * sin(y)
\end{verbatim}

\begin{Shaded}
\begin{Highlighting}[]
\FunctionTok{cat}\NormalTok{(}\StringTok{"∂²f/∂x∂y ="}\NormalTok{, }\FunctionTok{as.character}\NormalTok{(}\FunctionTok{body}\NormalTok{(fxy)), }\StringTok{"}\SpecialCharTok{\textbackslash{}n}\StringTok{"}\NormalTok{)}
\end{Highlighting}
\end{Shaded}

\begin{verbatim}
## ∂²f/∂x∂y = { .e1 <- x * y (1 + .e1) * exp(.e1) + 3 * (x^2 * cos(y))
\end{verbatim}

\section{Derivadas parciales de orden superior}\label{derivadas-parciales-de-orden-superior}

Las segundas derivadas parciales son:

\[
\begin{aligned}
f_{xx} &= \frac{\partial^2 f}{\partial x^2} = \frac{\partial}{\partial x}\left(\frac{\partial f}{\partial x}\right) \\
f_{yy} &= \frac{\partial^2 f}{\partial y^2} = \frac{\partial}{\partial y}\left(\frac{\partial f}{\partial y}\right) \\
f_{xy} &= \frac{\partial^2 f}{\partial y \partial x} = \frac{\partial}{\partial y}\left(\frac{\partial f}{\partial x}\right) \\
f_{yx} &= \frac{\partial^2 f}{\partial x \partial y} = \frac{\partial}{\partial x}\left(\frac{\partial f}{\partial y}\right)
\end{aligned}
\]

\subsection{Teorema de Clairaut}\label{teorema-de-clairaut}

Si \(f_{xy}\) y \(f_{yx}\) son continuas en una región, entonces:

\[
f_{xy} = f_{yx}
\]

Es decir, el orden de derivación no importa (bajo condiciones de regularidad).

🎧 Podcast: 03\_derivadas\_parciales\_teorema-de-clairaut

Tu navegador no soporta el elemento de audio.

\section{Plano tangente y aproximación lineal}\label{plano-tangente-y-aproximaciuxf3n-lineal}

\subsection{Ecuación del plano tangente}\label{ecuaciuxf3n-del-plano-tangente}

Si \(f\) es diferenciable en \((a, b)\), el \textbf{plano tangente} a la superficie \(z = f(x, y)\) en el punto \((a, b, f(a, b))\) tiene ecuación:

\[
z = f(a, b) + f_x(a, b)(x - a) + f_y(a, b)(y - b)
\]

\subsection{Aproximación lineal (diferenciabilidad)}\label{aproximaciuxf3n-lineal-diferenciabilidad}

La función \(f\) es \textbf{diferenciable} en \((a, b)\) si puede aproximarse linealmente cerca de ese punto:

\[
f(x, y) \approx f(a, b) + f_x(a, b)(x - a) + f_y(a, b)(y - b)
\]

El error de esta aproximación tiende a cero más rápidamente que la distancia \(\sqrt{(x-a)^2 + (y-b)^2}\).

\subsection{Ejemplo resuelto}\label{ejemplo-resuelto-10}

\begin{quote}
\textbf{Ejemplo 3.} Halle la ecuación del plano tangente a \(f(x, y) = x^2 + y^2\) en \((1, 2)\).
\end{quote}

\textbf{Solución.} \(f(1, 2) = 5\), \(f_x = 2x \implies f_x(1, 2) = 2\), \(f_y = 2y \implies f_y(1, 2) = 4\).

\[
z = 5 + 2(x - 1) + 4(y - 2) = 2x + 4y - 5
\]

\section{Regla de la cadena}\label{regla-de-la-cadena}

\subsection{Caso 1: Función compuesta de una variable}\label{caso-1-funciuxf3n-compuesta-de-una-variable}

Si \(z = f(x, y)\) donde \(x = x(t)\) e \(y = y(t)\), entonces:

\[
\frac{dz}{dt} = \frac{\partial f}{\partial x}\frac{dx}{dt} + \frac{\partial f}{\partial y}\frac{dy}{dt}
\]

\subsection{Caso 2: Función compuesta de dos variables}\label{caso-2-funciuxf3n-compuesta-de-dos-variables}

Si \(z = f(x, y)\) donde \(x = x(s, t)\) e \(y = y(s, t)\), entonces:

\[
\begin{aligned}
\frac{\partial z}{\partial s} &= \frac{\partial f}{\partial x}\frac{\partial x}{\partial s} + \frac{\partial f}{\partial y}\frac{\partial y}{\partial s} \\
\frac{\partial z}{\partial t} &= \frac{\partial f}{\partial x}\frac{\partial x}{\partial t} + \frac{\partial f}{\partial y}\frac{\partial y}{\partial t}
\end{aligned}
\]

\subsection{Ejemplo resuelto}\label{ejemplo-resuelto-11}

\begin{quote}
\textbf{Ejemplo 4.} Sea \(z = x^2 y\) donde \(x = \sin t\) e \(y = \cos t\). Calcule \(dz/dt\).
\end{quote}

\textbf{Solución.}

\[
\frac{dz}{dt} = (2xy)(\cos t) + (x^2)(-\sin t) = 2\sin t \cos^2 t - \sin^3 t
\]

\section{Derivada direccional y gradiente}\label{derivada-direccional-y-gradiente}

\subsection{Derivada direccional}\label{derivada-direccional}

La \textbf{derivada direccional} de \(f\) en \((a, b)\) en la dirección del vector unitario \(\vec{u} = (u_1, u_2)\) mide la tasa de cambio de \(f\) en esa dirección:

\[
D_{\vec{u}} f(a, b) = \lim_{h \to 0} \frac{f(a + hu_1, b + hu_2) - f(a, b)}{h}
\]

Si \(f\) es diferenciable, esta fórmula se simplifica a:

\[
D_{\vec{u}} f(a, b) = f_x(a, b) u_1 + f_y(a, b) u_2
\]

\subsection{Vector gradiente}\label{vector-gradiente}

El \textbf{gradiente} de \(f\) es el vector de las derivadas parciales:

\[
\nabla f = (f_x, f_y)
\]

En tres variables:

\[
\nabla f = (f_x, f_y, f_z)
\]

La derivada direccional puede escribirse como:

\[
D_{\vec{u}} f = \nabla f \cdot \vec{u}
\]

\subsection{Propiedades del gradiente}\label{propiedades-del-gradiente}

\begin{enumerate}
\def\labelenumi{\arabic{enumi}.}
\tightlist
\item
  \textbf{Dirección de máximo crecimiento}: el gradiente apunta en la dirección de máxima tasa de crecimiento de \(f\).
\item
  \textbf{Magnitud del máximo crecimiento}: \(\|\nabla f\|\) es la máxima tasa de crecimiento.
\item
  \textbf{Perpendicularidad}: el gradiente es perpendicular a las curvas (o superficies) de nivel.
\end{enumerate}

🎧 Podcast: 03\_derivadas\_parciales\_propiedades-del-gradiente

Tu navegador no soporta el elemento de audio.

\subsection{Ejemplo resuelto}\label{ejemplo-resuelto-12}

\begin{quote}
\textbf{Ejemplo 5.} Sea \(f(x, y) = x^2 + y^2\). Calcule el gradiente en \((1, 2)\), la derivada direccional en la dirección \(\vec{u} = (\frac{3}{5}, \frac{4}{5})\), y la dirección de máximo crecimiento.
\end{quote}

\textbf{Solución.} \(\nabla f = (2x, 2y)\), por tanto \(\nabla f(1, 2) = (2, 4)\).

Derivada direccional:

\[
D_{\vec{u}} f(1, 2) = (2, 4) \cdot \left(\frac{3}{5}, \frac{4}{5}\right) = \frac{6}{5} + \frac{16}{5} = \frac{22}{5}
\]

Dirección de máximo crecimiento: \(\nabla f(1, 2) = (2, 4)\), o normalizado: \(\frac{(2, 4)}{\sqrt{20}} = \frac{(1, 2)}{\sqrt{5}}\).

Máxima tasa de crecimiento: \(\|\nabla f(1, 2)\| = \sqrt{4 + 16} = \sqrt{20} = 2\sqrt{5}\).

\section{Visualización interactiva: plano tangente y gradiente}\label{visualizaciuxf3n-interactiva-plano-tangente-y-gradiente}

Explore el siguiente applet de GeoGebra para visualizar el plano tangente a una superficie y el vector gradiente en el plano \(xy\):

\section{Ejercicios propuestos}\label{ejercicios-propuestos-3}

\begin{enumerate}
\def\labelenumi{\arabic{enumi}.}
\tightlist
\item
  Calcule las derivadas parciales de \(f(x, y) = x^3 y^2 + \ln(xy)\).
\item
  Verifique el teorema de Clairaut para \(f(x, y) = x^3 y^2 + x\sin(y)\).
\item
  Halle la ecuación del plano tangente a \(z = e^{x-y}\) en \((1, 1)\).
\item
  Use la regla de la cadena para calcular \(\partial z / \partial s\) si \(z = x^2 + y^2\), \(x = s + t\), \(y = s - t\).
\item
  Calcule la derivada direccional de \(f(x, y) = x^2 y\) en \((2, 1)\) en la dirección del vector \((1, 1)\).
\item
  Halle la dirección de máximo crecimiento de \(f(x, y, z) = xyz\) en \((1, 2, 3)\).
\end{enumerate}

\chapter{Máximos, Mínimos y Optimización}\label{maximos-minimos}

🎧 Podcast del capítulo

Tu navegador no soporta el elemento de audio.

Una de las aplicaciones más importantes del cálculo diferencial es la búsqueda de valores extremos de funciones. En este capítulo se desarrollan los criterios para identificar máximos y mínimos locales de funciones de dos y tres variables, utilizando el Hessiano y los multiplicadores de Lagrange.

\section{Puntos críticos}\label{puntos-cruxedticos}

\subsection{Definición}\label{definiciuxf3n-4}

Un punto \((a, b)\) es un \textbf{punto crítico} de \(f(x, y)\) si:

\begin{enumerate}
\def\labelenumi{\arabic{enumi}.}
\tightlist
\item
  \(\nabla f(a, b) = \vec{0}\), es decir, \(f_x(a, b) = 0\) y \(f_y(a, b) = 0\); o
\item
  Alguna de las derivadas parciales no existe en \((a, b)\).
\end{enumerate}

Los puntos críticos de tipo 1 se llaman \textbf{puntos estacionarios}.

🎧 Podcast: 04\_maximos\_minimos\_definicion

Tu navegador no soporta el elemento de audio.

\subsection{Extremos locales}\label{extremos-locales}

\begin{itemize}
\tightlist
\item
  \(f(a, b)\) es un \textbf{máximo local} si \(f(x, y) \leq f(a, b)\) para todo \((x, y)\) en una vecindad de \((a, b)\).
\item
  \(f(a, b)\) es un \textbf{mínimo local} si \(f(x, y) \geq f(a, b)\) para todo \((x, y)\) en una vecindad de \((a, b)\).
\item
  Un punto crítico que no es ni máximo ni mínimo local se denomina \textbf{punto de silla}.
\end{itemize}

\section{Criterio del Hessiano para funciones de dos variables}\label{criterio-del-hessiano-para-funciones-de-dos-variables}

\subsection{Matriz Hessiana}\label{matriz-hessiana}

La \textbf{matriz Hessiana} de \(f\) en \((a, b)\) es:

\[
H(f) = \begin{pmatrix}
f_{xx} & f_{xy} \\
f_{yx} & f_{yy}
\end{pmatrix}
\]

\subsection{Discriminante}\label{discriminante}

El \textbf{discriminante} (o determinante del Hessiano) es:

\[
D = f_{xx} f_{yy} - (f_{xy})^2 = \det(H)
\]

\subsection{Criterio de clasificación}\label{criterio-de-clasificaciuxf3n}

Sea \((a, b)\) un punto crítico de \(f\) con segundas derivadas parciales continuas en una vecindad de \((a, b)\):

\begin{longtable}[]{@{}ll@{}}
\toprule\noalign{}
Condición & Clasificación \\
\midrule\noalign{}
\endhead
\bottomrule\noalign{}
\endlastfoot
\(D > 0\) y \(f_{xx} > 0\) & \textbf{Mínimo local} \\
\(D > 0\) y \(f_{xx} < 0\) & \textbf{Máximo local} \\
\(D < 0\) & \textbf{Punto de silla} \\
\(D = 0\) & \textbf{Criterio inconcluso} \\
\end{longtable}

\subsection{Ejemplo resuelto}\label{ejemplo-resuelto-13}

\begin{quote}
\textbf{Ejemplo 1.} Clasifique los puntos críticos de \(f(x, y) = x^3 + y^3 - 3xy\).
\end{quote}

\textbf{Solución.} Las derivadas parciales son:

\[
f_x = 3x^2 - 3y = 0 \implies y = x^2
\]

\[
f_y = 3y^2 - 3x = 0 \implies x = y^2
\]

Sustituyendo \(y = x^2\) en la segunda ecuación: \(x = x^4\), es decir, \(x(x^3 - 1) = 0\).

\begin{itemize}
\tightlist
\item
  Si \(x = 0\), entonces \(y = 0\). Punto crítico: \((0, 0)\).
\item
  Si \(x = 1\), entonces \(y = 1\). Punto crítico: \((1, 1)\).
\end{itemize}

Las segundas derivadas son \(f_{xx} = 6x\), \(f_{yy} = 6y\), \(f_{xy} = -3\).

\textbf{En \((0, 0)\):} \(f_{xx} = 0\), \(f_{yy} = 0\), \(f_{xy} = -3\).

\[
D = (0)(0) - (-3)^2 = -9 < 0 \implies \text{Punto de silla}
\]

\textbf{En \((1, 1)\):} \(f_{xx} = 6\), \(f_{yy} = 6\), \(f_{xy} = -3\).

\[
D = (6)(6) - (-3)^2 = 36 - 9 = 27 > 0, \quad f_{xx} = 6 > 0 \implies \text{Mínimo local}
\]

\section{Herramienta de cálculo: Hessiano automático}\label{herramienta-de-cuxe1lculo-hessiano-automuxe1tico}

A continuación se presenta una herramienta de cálculo simbólico que calcula automáticamente el Hessiano, lo evalúa en un punto crítico y clasifica el punto. El usuario puede modificar la función y el punto en el código para analizar sus propios ejemplos.

\subsection{\texorpdfstring{Ejemplo: \(f(x, y) = x^3 + y^3 - 3xy\) en \((1, 1)\)}{Ejemplo: f(x, y) = x\^{}3 + y\^{}3 - 3xy en (1, 1)}}\label{ejemplo-fx-y-x3-y3---3xy-en-1-1}

\begin{Shaded}
\begin{Highlighting}[]
\FunctionTok{library}\NormalTok{(Deriv)}

\CommentTok{\# Defina aquí su función y el punto crítico}
\NormalTok{f }\OtherTok{\textless{}{-}} \ControlFlowTok{function}\NormalTok{(x, y) x}\SpecialCharTok{\^{}}\DecValTok{3} \SpecialCharTok{+}\NormalTok{ y}\SpecialCharTok{\^{}}\DecValTok{3} \SpecialCharTok{{-}} \DecValTok{3}\SpecialCharTok{*}\NormalTok{x}\SpecialCharTok{*}\NormalTok{y}
\NormalTok{a }\OtherTok{\textless{}{-}} \DecValTok{1}  \CommentTok{\# coordenada x del punto crítico}
\NormalTok{b }\OtherTok{\textless{}{-}} \DecValTok{1}  \CommentTok{\# coordenada y del punto crítico}

\CommentTok{\# Cálculo simbólico de derivadas}
\NormalTok{fx }\OtherTok{\textless{}{-}} \FunctionTok{Deriv}\NormalTok{(f, }\StringTok{"x"}\NormalTok{)}
\NormalTok{fy }\OtherTok{\textless{}{-}} \FunctionTok{Deriv}\NormalTok{(f, }\StringTok{"y"}\NormalTok{)}
\NormalTok{fxx }\OtherTok{\textless{}{-}} \FunctionTok{Deriv}\NormalTok{(fx, }\StringTok{"x"}\NormalTok{)}
\NormalTok{fyy }\OtherTok{\textless{}{-}} \FunctionTok{Deriv}\NormalTok{(fy, }\StringTok{"y"}\NormalTok{)}
\NormalTok{fxy }\OtherTok{\textless{}{-}} \FunctionTok{Deriv}\NormalTok{(fx, }\StringTok{"y"}\NormalTok{)}

\CommentTok{\# Evaluación en el punto (a, b)}
\NormalTok{H11 }\OtherTok{\textless{}{-}} \FunctionTok{fxx}\NormalTok{(a, b)}
\NormalTok{H22 }\OtherTok{\textless{}{-}} \FunctionTok{fyy}\NormalTok{(a, b)}
\NormalTok{H12 }\OtherTok{\textless{}{-}} \FunctionTok{fxy}\NormalTok{(a, b)}
\NormalTok{D }\OtherTok{\textless{}{-}}\NormalTok{ H11 }\SpecialCharTok{*}\NormalTok{ H22 }\SpecialCharTok{{-}}\NormalTok{ H12}\SpecialCharTok{\^{}}\DecValTok{2}

\CommentTok{\# Clasificación}
\ControlFlowTok{if}\NormalTok{ (D }\SpecialCharTok{\textgreater{}} \DecValTok{0}\NormalTok{) \{}
  \ControlFlowTok{if}\NormalTok{ (H11 }\SpecialCharTok{\textgreater{}} \DecValTok{0}\NormalTok{) clasif }\OtherTok{\textless{}{-}} \StringTok{"MÍNIMO LOCAL"}
  \ControlFlowTok{else} \ControlFlowTok{if}\NormalTok{ (H11 }\SpecialCharTok{\textless{}} \DecValTok{0}\NormalTok{) clasif }\OtherTok{\textless{}{-}} \StringTok{"MÁXIMO LOCAL"}
  \ControlFlowTok{else}\NormalTok{ clasif }\OtherTok{\textless{}{-}} \StringTok{"Indeterminado (D \textgreater{} 0 pero f\_xx = 0)"}
\NormalTok{\} }\ControlFlowTok{else} \ControlFlowTok{if}\NormalTok{ (D }\SpecialCharTok{\textless{}} \DecValTok{0}\NormalTok{) \{}
\NormalTok{  clasif }\OtherTok{\textless{}{-}} \StringTok{"PUNTO DE SILLA"}
\NormalTok{\} }\ControlFlowTok{else}\NormalTok{ \{}
\NormalTok{  clasif }\OtherTok{\textless{}{-}} \StringTok{"CRITERIO INCONCLUSO (D = 0)"}
\NormalTok{\}}

\FunctionTok{cat}\NormalTok{(}\StringTok{"=== ANÁLISIS DEL HESSIANO ===}\SpecialCharTok{\textbackslash{}n\textbackslash{}n}\StringTok{"}\NormalTok{)}
\end{Highlighting}
\end{Shaded}

\begin{verbatim}
## === ANÁLISIS DEL HESSIANO ===
\end{verbatim}

\begin{Shaded}
\begin{Highlighting}[]
\FunctionTok{cat}\NormalTok{(}\StringTok{"Función: f(x,y) = x³ + y³ {-} 3xy}\SpecialCharTok{\textbackslash{}n}\StringTok{"}\NormalTok{)}
\end{Highlighting}
\end{Shaded}

\begin{verbatim}
## Función: f(x,y) = x³ + y³ - 3xy
\end{verbatim}

\begin{Shaded}
\begin{Highlighting}[]
\FunctionTok{cat}\NormalTok{(}\StringTok{"Punto crítico: ("}\NormalTok{, a, }\StringTok{", "}\NormalTok{, b, }\StringTok{")}\SpecialCharTok{\textbackslash{}n\textbackslash{}n}\StringTok{"}\NormalTok{, }\AttributeTok{sep =} \StringTok{""}\NormalTok{)}
\end{Highlighting}
\end{Shaded}

\begin{verbatim}
## Punto crítico: (1, 1)
\end{verbatim}

\begin{Shaded}
\begin{Highlighting}[]
\FunctionTok{cat}\NormalTok{(}\StringTok{"f\_xx ="}\NormalTok{, H11, }\StringTok{"}\SpecialCharTok{\textbackslash{}n}\StringTok{"}\NormalTok{)}
\end{Highlighting}
\end{Shaded}

\begin{verbatim}
## f_xx = 6
\end{verbatim}

\begin{Shaded}
\begin{Highlighting}[]
\FunctionTok{cat}\NormalTok{(}\StringTok{"f\_yy ="}\NormalTok{, H22, }\StringTok{"}\SpecialCharTok{\textbackslash{}n}\StringTok{"}\NormalTok{)}
\end{Highlighting}
\end{Shaded}

\begin{verbatim}
## f_yy = 6
\end{verbatim}

\begin{Shaded}
\begin{Highlighting}[]
\FunctionTok{cat}\NormalTok{(}\StringTok{"f\_xy ="}\NormalTok{, H12, }\StringTok{"}\SpecialCharTok{\textbackslash{}n\textbackslash{}n}\StringTok{"}\NormalTok{)}
\end{Highlighting}
\end{Shaded}

\begin{verbatim}
## f_xy = -3
\end{verbatim}

\begin{Shaded}
\begin{Highlighting}[]
\FunctionTok{cat}\NormalTok{(}\StringTok{"Discriminante D = f\_xx·f\_yy {-} (f\_xy)² ="}\NormalTok{, D, }\StringTok{"}\SpecialCharTok{\textbackslash{}n\textbackslash{}n}\StringTok{"}\NormalTok{)}
\end{Highlighting}
\end{Shaded}

\begin{verbatim}
## Discriminante D = f_xx·f_yy - (f_xy)² = 27
\end{verbatim}

\begin{Shaded}
\begin{Highlighting}[]
\FunctionTok{cat}\NormalTok{(}\StringTok{"CLASIFICACIÓN:"}\NormalTok{, clasif, }\StringTok{"}\SpecialCharTok{\textbackslash{}n}\StringTok{"}\NormalTok{)}
\end{Highlighting}
\end{Shaded}

\begin{verbatim}
## CLASIFICACIÓN: MÍNIMO LOCAL
\end{verbatim}

\subsection{\texorpdfstring{Segundo ejemplo: \(f(x, y) = x^2 + y^2\) en \((0, 0)\)}{Segundo ejemplo: f(x, y) = x\^{}2 + y\^{}2 en (0, 0)}}\label{segundo-ejemplo-fx-y-x2-y2-en-0-0}

\begin{Shaded}
\begin{Highlighting}[]
\NormalTok{f2 }\OtherTok{\textless{}{-}} \ControlFlowTok{function}\NormalTok{(x, y) x}\SpecialCharTok{\^{}}\DecValTok{2} \SpecialCharTok{+}\NormalTok{ y}\SpecialCharTok{\^{}}\DecValTok{2}
\NormalTok{a2 }\OtherTok{\textless{}{-}} \DecValTok{0}\NormalTok{; b2 }\OtherTok{\textless{}{-}} \DecValTok{0}

\NormalTok{fx2 }\OtherTok{\textless{}{-}} \FunctionTok{Deriv}\NormalTok{(f2, }\StringTok{"x"}\NormalTok{)}
\NormalTok{fy2 }\OtherTok{\textless{}{-}} \FunctionTok{Deriv}\NormalTok{(f2, }\StringTok{"y"}\NormalTok{)}
\NormalTok{fxx2 }\OtherTok{\textless{}{-}} \FunctionTok{Deriv}\NormalTok{(fx2, }\StringTok{"x"}\NormalTok{)}
\NormalTok{fyy2 }\OtherTok{\textless{}{-}} \FunctionTok{Deriv}\NormalTok{(fy2, }\StringTok{"y"}\NormalTok{)}
\NormalTok{fxy2 }\OtherTok{\textless{}{-}} \FunctionTok{Deriv}\NormalTok{(fx2, }\StringTok{"y"}\NormalTok{)}

\NormalTok{H11\_2 }\OtherTok{\textless{}{-}} \FunctionTok{fxx2}\NormalTok{(a2, b2)}
\NormalTok{H22\_2 }\OtherTok{\textless{}{-}} \FunctionTok{fyy2}\NormalTok{(a2, b2)}
\NormalTok{H12\_2 }\OtherTok{\textless{}{-}} \FunctionTok{fxy2}\NormalTok{(a2, b2)}
\NormalTok{D2 }\OtherTok{\textless{}{-}}\NormalTok{ H11\_2 }\SpecialCharTok{*}\NormalTok{ H22\_2 }\SpecialCharTok{{-}}\NormalTok{ H12\_2}\SpecialCharTok{\^{}}\DecValTok{2}

\ControlFlowTok{if}\NormalTok{ (D2 }\SpecialCharTok{\textgreater{}} \DecValTok{0} \SpecialCharTok{\&\&}\NormalTok{ H11\_2 }\SpecialCharTok{\textgreater{}} \DecValTok{0}\NormalTok{) clasif2 }\OtherTok{\textless{}{-}} \StringTok{"MÍNIMO LOCAL"}
\FunctionTok{cat}\NormalTok{(}\StringTok{"f(x,y) = x² + y² en (0,0): D ="}\NormalTok{, D2, }\StringTok{", f\_xx ="}\NormalTok{, H11\_2, }\StringTok{"→"}\NormalTok{, clasif2, }\StringTok{"}\SpecialCharTok{\textbackslash{}n}\StringTok{"}\NormalTok{)}
\end{Highlighting}
\end{Shaded}

\begin{verbatim}
## f(x,y) = x² + y² en (0,0): D = 4 , f_xx = 2 → MÍNIMO LOCAL
\end{verbatim}

\section{Extremos absolutos en regiones compactas}\label{extremos-absolutos-en-regiones-compactas}

El \textbf{teorema de Weierstrass} garantiza que una función continua sobre un conjunto cerrado y acotado (compacto) alcanza su máximo y mínimo absolutos.

\subsection{Método para encontrar extremos absolutos}\label{muxe9todo-para-encontrar-extremos-absolutos}

\begin{enumerate}
\def\labelenumi{\arabic{enumi}.}
\tightlist
\item
  Encuentre los puntos críticos en el interior de la región.
\item
  Encuentre los extremos de \(f\) restringida a la frontera (usando parametrización o multiplicadores de Lagrange).
\item
  Evalúe \(f\) en todos los puntos candidatos y compare.
\end{enumerate}

\subsection{Ejemplo resuelto}\label{ejemplo-resuelto-14}

\begin{quote}
\textbf{Ejemplo 2.} Halle los extremos absolutos de \(f(x, y) = x^2 + y^2 - x - y\) en el disco \(x^2 + y^2 \leq 1\).
\end{quote}

\textbf{Solución.} \textbf{Interior:} \(f_x = 2x - 1 = 0\), \(f_y = 2y - 1 = 0\). Punto crítico: \((\frac{1}{2}, \frac{1}{2})\).

\(f(\frac{1}{2}, \frac{1}{2}) = \frac{1}{4} + \frac{1}{4} - \frac{1}{2} - \frac{1}{2} = -\frac{1}{2}\).

\textbf{Frontera:} Parametrizamos \(x = \cos t\), \(y = \sin t\). Entonces:

\[
g(t) = 1 - \cos t - \sin t
\]

\(g'(t) = \sin t - \cos t = 0 \implies \tan t = 1\), es decir, \(t = \pi/4\) o \(t = 5\pi/4\).

\begin{itemize}
\tightlist
\item
  \(t = \pi/4\): \((\frac{\sqrt{2}}{2}, \frac{\sqrt{2}}{2})\), \(f = 1 - \sqrt{2} \approx -0.414\)
\item
  \(t = 5\pi/4\): \((-\frac{\sqrt{2}}{2}, -\frac{\sqrt{2}}{2})\), \(f = 1 + \sqrt{2} \approx 2.414\)
\end{itemize}

Comparando: \textbf{mínimo absoluto} \(-\frac{1}{2}\) en \((\frac{1}{2}, \frac{1}{2})\), \textbf{máximo absoluto} \(1 + \sqrt{2}\) en \((-\frac{\sqrt{2}}{2}, -\frac{\sqrt{2}}{2})\).

\section{Multiplicadores de Lagrange}\label{multiplicadores-de-lagrange}

El método de los multiplicadores de Lagrange permite encontrar extremos de una función \(f(x, y)\) sujeta a una restricción \(g(x, y) = 0\).

\subsection{Condición necesaria}\label{condiciuxf3n-necesaria}

Si \(f\) tiene un extremo en \((a, b)\) sujeto a \(g(x, y) = 0\), y \(\nabla g(a, b) \neq \vec{0}\), entonces existe un escalar \(\lambda\) (llamado \textbf{multiplicador de Lagrange}) tal que:

\[
\nabla f(a, b) = \lambda \nabla g(a, b)
\]

Esto, junto con la restricción \(g(a, b) = 0\), forma un sistema de ecuaciones.

\subsection{Función de Lagrange}\label{funciuxf3n-de-lagrange}

Definimos:

\[
\mathcal{L}(x, y, \lambda) = f(x, y) - \lambda g(x, y)
\]

Los puntos críticos de \(\mathcal{L}\) son los candidatos a extremos condicionados.

\subsection{Ejemplo resuelto}\label{ejemplo-resuelto-15}

\begin{quote}
\textbf{Ejemplo 3.} Halle el punto de la recta \(x + y = 1\) más cercano al origen.
\end{quote}

\textbf{Solución.} Queremos minimizar \(f(x, y) = x^2 + y^2\) sujeto a \(g(x, y) = x + y - 1 = 0\).

La función de Lagrange es:

\[
\mathcal{L} = x^2 + y^2 - \lambda(x + y - 1)
\]

Las condiciones son:

\[
\begin{aligned}
\mathcal{L}_x &= 2x - \lambda = 0 \\
\mathcal{L}_y &= 2y - \lambda = 0 \\
\mathcal{L}_\lambda &= -(x + y - 1) = 0
\end{aligned}
\]

De las primeras dos ecuaciones: \(2x = 2y = \lambda \implies x = y\). Sustituyendo en la tercera: \(2x = 1 \implies x = y = \frac{1}{2}\).

El punto más cercano es \((\frac{1}{2}, \frac{1}{2})\) y la distancia mínima es \(\sqrt{\frac{1}{4} + \frac{1}{4}} = \frac{\sqrt{2}}{2}\).

\section{Extensión a tres variables}\label{extensiuxf3n-a-tres-variables}

Para una función \(f(x, y, z)\) sujeta a una restricción \(g(x, y, z) = 0\), la condición es análoga:

\[
\nabla f = \lambda \nabla g, \qquad g = 0
\]

Para dos restricciones \(g = 0\) y \(h = 0\):

\[
\nabla f = \lambda \nabla g + \mu \nabla h, \qquad g = 0, \quad h = 0
\]

\section{Visualización interactiva: puntos críticos y sillas de montar}\label{visualizaciuxf3n-interactiva-puntos-cruxedticos-y-sillas-de-montar}

Explore el siguiente applet de GeoGebra para visualizar máximos, mínimos y puntos de silla en superficies:

\section{Ejercicios propuestos}\label{ejercicios-propuestos-4}

\begin{enumerate}
\def\labelenumi{\arabic{enumi}.}
\tightlist
\item
  Clasifique los puntos críticos de \(f(x, y) = x^3 - 3x + y^2\).
\item
  Halle y clasifique los puntos críticos de \(f(x, y) = x^4 + y^4 - 4xy\).
\item
  Determine los extremos absolutos de \(f(x, y) = xy\) en el cuadrado \([0, 1] \times [0, 1]\).
\item
  Use multiplicadores de Lagrange para hallar el rectángulo de área máxima con perímetro fijo \(P\).
\item
  Halle los puntos de la esfera \(x^2 + y^2 + z^2 = 1\) donde \(f(x, y, z) = x + 2y + 3z\) alcanza sus extremos.
\item
  Demuestre que de todos los paralelepípedos rectangulares de volumen fijo \(V\), el cubo tiene la menor superficie.
\end{enumerate}

\chapter{Superficies Cuádricas}\label{cuadricas}

🎧 Podcast del capítulo

Tu navegador no soporta el elemento de audio.

Las superficies cuádricas constituyen la familia más importante de superficies en geometría analítica. Son el análogo tridimensional de las secciones cónicas y aparecen de manera natural en física, ingeniería, arquitectura y gráficos por computadora.

\section{Ecuación general y clasificación}\label{ecuaciuxf3n-general-y-clasificaciuxf3n}

Una \textbf{superficie cuádrica} es el conjunto de puntos \((x, y, z) \in \mathbb{R}^3\) que satisfacen una ecuación algebraica de segundo grado:

\[
Ax^2 + By^2 + Cz^2 + Dxy + Exz + Fyz + Gx + Hy + Iz + J = 0
\]

donde \(A, B, \dots, J\) son constantes reales y al menos uno de los coeficientes cuadráticos es distinto de cero.

Mediante rotaciones y traslaciones de coordenadas, toda ecuación cuádrica puede reducirse a una de las \textbf{formas canónicas} que estudiaremos en este capítulo. Cuando la superficie tiene su centro en el origen y sus ejes de simetría coinciden con los ejes coordenados, la ecuación no contiene términos cruzados (\(D = E = F = 0\)) ni términos lineales (\(G = H = I = 0\)).

\subsection{Estrategia de análisis}\label{estrategia-de-anuxe1lisis}

Para cada superficie cuádrica estudiaremos:

\begin{enumerate}
\def\labelenumi{\arabic{enumi}.}
\tightlist
\item
  Su ecuación canónica.
\item
  Las \textbf{trazas} en los planos coordenados y en planos paralelos a ellos.
\item
  Sus simetrías.
\item
  Si está acotada o se extiende al infinito.
\item
  Propiedades geométricas destacadas (regladas, puntos de silla, etc.).
\item
  Un \textbf{ejemplo resuelto}.
\item
  Una \textbf{visualización interactiva}.
\end{enumerate}

\begin{center}\rule{0.5\linewidth}{0.5pt}\end{center}

\section{Elipsoide}\label{elipsoide}

\subsection{Ecuación canónica}\label{ecuaciuxf3n-canuxf3nica}

\[
\frac{x^2}{a^2} + \frac{y^2}{b^2} + \frac{z^2}{c^2} = 1
\]

donde \(a, b, c > 0\) son las longitudes de los semiejes en las direcciones \(x\), \(y\) y \(z\) respectivamente.

\subsection{Caso particular: la esfera}\label{caso-particular-la-esfera}

Cuando \(a = b = c = r\), la ecuación se reduce a:

\[
x^2 + y^2 + z^2 = r^2
\]

que es una esfera de radio \(r\) centrada en el origen.

\subsection{Trazas}\label{trazas}

\begin{longtable}[]{@{}
  >{\raggedright\arraybackslash}p{(\columnwidth - 4\tabcolsep) * \real{0.2000}}
  >{\raggedright\arraybackslash}p{(\columnwidth - 4\tabcolsep) * \real{0.6286}}
  >{\raggedright\arraybackslash}p{(\columnwidth - 4\tabcolsep) * \real{0.1714}}@{}}
\toprule\noalign{}
\begin{minipage}[b]{\linewidth}\raggedright
Plano
\end{minipage} & \begin{minipage}[b]{\linewidth}\raggedright
Ecuación de la traza
\end{minipage} & \begin{minipage}[b]{\linewidth}\raggedright
Tipo
\end{minipage} \\
\midrule\noalign{}
\endhead
\bottomrule\noalign{}
\endlastfoot
\(xy\) (\(z = 0\)) & \(\frac{x^2}{a^2} + \frac{y^2}{b^2} = 1\) & Elipse \\
\(xz\) (\(y = 0\)) & \(\frac{x^2}{a^2} + \frac{z^2}{c^2} = 1\) & Elipse \\
\(yz\) (\(x = 0\)) & \(\frac{y^2}{b^2} + \frac{z^2}{c^2} = 1\) & Elipse \\
\(z = k\) (\(|k| < c\)) & \(\frac{x^2}{a^2} + \frac{y^2}{b^2} = 1 - \frac{k^2}{c^2}\) & Elipse \\
\(z = k\) (\(|k| = c\)) & Punto \((0, 0, k)\) & Punto \\
\(z = k\) (\(|k| > c\)) & No hay intersección & Vacío \\
\end{longtable}

\subsection{Propiedades geométricas}\label{propiedades-geomuxe9tricas-1}

\begin{itemize}
\tightlist
\item
  Es una superficie \textbf{acotada} y \textbf{cerrada}.
\item
  Es simétrica respecto a todos los planos coordenados, ejes coordenados y el origen.
\item
  Los puntos \((\pm a, 0, 0)\), \((0, \pm b, 0)\) y \((0, 0, \pm c)\) son los \textbf{vértices} del elipsoide.
\item
  Si \(a > b > c\), el eje mayor es el \(x\), el medio el \(y\) y el menor el \(z\).
\end{itemize}

🎧 Podcast: 05\_cuadricas\_propiedades-geometricas

Tu navegador no soporta el elemento de audio.

\subsection{Ejemplo resuelto}\label{ejemplo-resuelto-16}

\begin{quote}
\textbf{Ejemplo 1.} Identifique y grafique la superficie \(4x^2 + 9y^2 + 36z^2 = 36\).
\end{quote}

\textbf{Solución.} Dividiendo entre 36:

\[
\frac{x^2}{9} + \frac{y^2}{4} + \frac{z^2}{1} = 1
\]

Se trata de un \textbf{elipsoide} con semiejes \(a = 3\), \(b = 2\), \(c = 1\). Las trazas en los planos coordenados son elipses de semiejes \((3, 2)\), \((3, 1)\) y \((2, 1)\) respectivamente. La superficie es acotada y se extiende desde \(x = -3\) hasta \(x = 3\), desde \(y = -2\) hasta \(y = 2\), y desde \(z = -1\) hasta \(z = 1\).

\label{fig:elipsoide-ejemplo}Elipsoide del Ejemplo 1 con semiejes \(a = 3\), \(b = 2\), \(c = 1\).

\begin{center}\rule{0.5\linewidth}{0.5pt}\end{center}

\section{Hiperboloide de una hoja}\label{hiperboloide-de-una-hoja}

\subsection{Ecuación canónica}\label{ecuaciuxf3n-canuxf3nica-1}

\[
\frac{x^2}{a^2} + \frac{y^2}{b^2} - \frac{z^2}{c^2} = 1
\]

Observe el patrón de signos: dos positivos y uno negativo, igual a 1.

\subsection{Trazas}\label{trazas-1}

\begin{longtable}[]{@{}
  >{\raggedright\arraybackslash}p{(\columnwidth - 4\tabcolsep) * \real{0.2000}}
  >{\raggedright\arraybackslash}p{(\columnwidth - 4\tabcolsep) * \real{0.6286}}
  >{\raggedright\arraybackslash}p{(\columnwidth - 4\tabcolsep) * \real{0.1714}}@{}}
\toprule\noalign{}
\begin{minipage}[b]{\linewidth}\raggedright
Plano
\end{minipage} & \begin{minipage}[b]{\linewidth}\raggedright
Ecuación de la traza
\end{minipage} & \begin{minipage}[b]{\linewidth}\raggedright
Tipo
\end{minipage} \\
\midrule\noalign{}
\endhead
\bottomrule\noalign{}
\endlastfoot
\(xy\) (\(z = 0\)) & \(\frac{x^2}{a^2} + \frac{y^2}{b^2} = 1\) & Elipse \\
\(xz\) (\(y = 0\)) & \(\frac{x^2}{a^2} - \frac{z^2}{c^2} = 1\) & Hipérbola \\
\(yz\) (\(x = 0\)) & \(\frac{y^2}{b^2} - \frac{z^2}{c^2} = 1\) & Hipérbola \\
\(z = k\) & \(\frac{x^2}{a^2} + \frac{y^2}{b^2} = 1 + \frac{k^2}{c^2}\) & Elipse \\
\end{longtable}

\subsection{Propiedades geométricas}\label{propiedades-geomuxe9tricas-2}

\begin{itemize}
\tightlist
\item
  Es una superficie \textbf{no acotada} que se extiende infinitamente en la dirección del eje asociado al término negativo (en este caso, el eje \(z\)).
\item
  Es \textbf{conexa}: consta de una sola pieza.
\item
  Es una superficie \textbf{reglada}: por cada uno de sus puntos pasan dos rectas completamente contenidas en la superficie. Esta propiedad tiene importantes aplicaciones en arquitectura y diseño estructural.
\item
  Es simétrica respecto a todos los planos coordenados, ejes coordenados y el origen.
\item
  Las trazas horizontales son siempre elipses, y las trazas verticales son hipérbolas.
\end{itemize}

🎧 Podcast: 05\_cuadricas\_propiedades-geometricas\_1

Tu navegador no soporta el elemento de audio.

\subsection{Ejemplo resuelto}\label{ejemplo-resuelto-17}

\begin{quote}
\textbf{Ejemplo 2.} Identifique la superficie \(x^2 + y^2 - \frac{z^2}{4} = 1\) y analice sus trazas.
\end{quote}

\textbf{Solución.} La ecuación corresponde a un \textbf{hiperboloide de una hoja} con \(a = b = 1\) y \(c = 2\). Como \(a = b\), las trazas horizontales son círculos; por esta razón se denomina \textbf{hiperboloide de una hoja de revolución} (o circular).

\begin{itemize}
\tightlist
\item
  Traza en \(z = 0\): \(x^2 + y^2 = 1\) (círculo de radio 1).
\item
  Traza en \(z = 2\): \(x^2 + y^2 = 2\) (círculo de radio \(\sqrt{2}\)).
\item
  Traza en \(y = 0\): \(x^2 - \frac{z^2}{4} = 1\) (hipérbola).
\end{itemize}

Observe que a medida que \(|z|\) aumenta, los círculos horizontales crecen de radio, lo cual explica la forma de ``copas'' de la superficie.

\label{fig:hiperboloide1-ejemplo}Hiperboloide de una hoja del Ejemplo 2.

\begin{center}\rule{0.5\linewidth}{0.5pt}\end{center}

\section{Hiperboloide de dos hojas}\label{hiperboloide-de-dos-hojas}

\subsection{Ecuación canónica}\label{ecuaciuxf3n-canuxf3nica-2}

\[
\frac{x^2}{a^2} + \frac{y^2}{b^2} - \frac{z^2}{c^2} = -1
\]

o de manera equivalente:

\[
\frac{z^2}{c^2} - \frac{x^2}{a^2} - \frac{y^2}{b^2} = 1
\]

El patrón de signos es: dos negativos y uno positivo, igual a 1 (o dos positivos y uno negativo, igual a \(-1\)).

\subsection{Trazas}\label{trazas-2}

\begin{longtable}[]{@{}
  >{\raggedright\arraybackslash}p{(\columnwidth - 4\tabcolsep) * \real{0.2000}}
  >{\raggedright\arraybackslash}p{(\columnwidth - 4\tabcolsep) * \real{0.6286}}
  >{\raggedright\arraybackslash}p{(\columnwidth - 4\tabcolsep) * \real{0.1714}}@{}}
\toprule\noalign{}
\begin{minipage}[b]{\linewidth}\raggedright
Plano
\end{minipage} & \begin{minipage}[b]{\linewidth}\raggedright
Ecuación de la traza
\end{minipage} & \begin{minipage}[b]{\linewidth}\raggedright
Tipo
\end{minipage} \\
\midrule\noalign{}
\endhead
\bottomrule\noalign{}
\endlastfoot
\(xy\) (\(z = 0\)) & \(\frac{x^2}{a^2} + \frac{y^2}{b^2} = -1\) & Vacío (no hay intersección) \\
\(xz\) (\(y = 0\)) & \(\frac{z^2}{c^2} - \frac{x^2}{a^2} = 1\) & Hipérbola \\
\(yz\) (\(x = 0\)) & \(\frac{z^2}{c^2} - \frac{y^2}{b^2} = 1\) & Hipérbola \\
\(z = k\) (\(|k| > c\)) & \(\frac{x^2}{a^2} + \frac{y^2}{b^2} = \frac{k^2}{c^2} - 1\) & Elipse \\
\(z = k\) (\(|k| = c\)) & Punto \((0, 0, k)\) & Punto \\
\(z = k\) (\(|k| < c\)) & No hay intersección & Vacío \\
\end{longtable}

\subsection{Propiedades geométricas}\label{propiedades-geomuxe9tricas-3}

\begin{itemize}
\tightlist
\item
  Es una superficie \textbf{no acotada}.
\item
  Consta de \textbf{dos hojas separadas}, una con \(z \geq c\) y otra con \(z \leq -c\). No es conexa.
\item
  \textbf{No es una superficie reglada}: no contiene rectas.
\item
  Es simétrica respecto a todos los planos coordenados, ejes coordenados y el origen.
\item
  Los vértices de las dos hojas están en \((0, 0, c)\) y \((0, 0, -c)\).
\end{itemize}

🎧 Podcast: 05\_cuadricas\_propiedades-geometricas\_2

Tu navegador no soporta el elemento de audio.

\subsection{Ejemplo resuelto}\label{ejemplo-resuelto-18}

\begin{quote}
\textbf{Ejemplo 3.} Analice la superficie \(-\frac{x^2}{4} - \frac{y^2}{9} + z^2 = 1\).
\end{quote}

\textbf{Solución.} Reescribiendo:

\[
z^2 - \frac{x^2}{4} - \frac{y^2}{9} = 1
\]

Se trata de un \textbf{hiperboloide de dos hojas} con \(a = 2\), \(b = 3\), \(c = 1\). Las dos hojas abren hacia arriba y hacia abajo a lo largo del eje \(z\). No hay superficie entre \(z = -1\) y \(z = 1\). Las trazas horizontales con \(|z| > 1\) son elipses cuyo tamaño aumenta a medida que \(|z|\) crece.

\label{fig:hiperboloide2-ejemplo}Hiperboloide de dos hojas del Ejemplo 3.

\begin{center}\rule{0.5\linewidth}{0.5pt}\end{center}

\section{Paraboloide elíptico}\label{paraboloide-eluxedptico}

\subsection{Ecuación canónica}\label{ecuaciuxf3n-canuxf3nica-3}

\[
\frac{x^2}{a^2} + \frac{y^2}{b^2} = \frac{z}{c}
\]

Si \(c > 0\), la superficie abre hacia arriba; si \(c < 0\), abre hacia abajo.

\subsection{Trazas}\label{trazas-3}

\begin{longtable}[]{@{}
  >{\raggedright\arraybackslash}p{(\columnwidth - 4\tabcolsep) * \real{0.2000}}
  >{\raggedright\arraybackslash}p{(\columnwidth - 4\tabcolsep) * \real{0.6286}}
  >{\raggedright\arraybackslash}p{(\columnwidth - 4\tabcolsep) * \real{0.1714}}@{}}
\toprule\noalign{}
\begin{minipage}[b]{\linewidth}\raggedright
Plano
\end{minipage} & \begin{minipage}[b]{\linewidth}\raggedright
Ecuación de la traza
\end{minipage} & \begin{minipage}[b]{\linewidth}\raggedright
Tipo
\end{minipage} \\
\midrule\noalign{}
\endhead
\bottomrule\noalign{}
\endlastfoot
\(xy\) (\(z = 0\)) & \(\frac{x^2}{a^2} + \frac{y^2}{b^2} = 0\) & Punto \((0, 0, 0)\) \\
\(xz\) (\(y = 0\)) & \(z = \frac{c}{a^2} x^2\) & Parábola \\
\(yz\) (\(x = 0\)) & \(z = \frac{c}{b^2} y^2\) & Parábola \\
\(z = k\) (\(k/c > 0\)) & \(\frac{x^2}{a^2} + \frac{y^2}{b^2} = \frac{k}{c}\) & Elipse \\
\end{longtable}

\subsection{Propiedades geométricas}\label{propiedades-geomuxe9tricas-4}

\begin{itemize}
\tightlist
\item
  Es una superficie \textbf{no acotada}.
\item
  Tiene un único \textbf{vértice} en el origen \((0, 0, 0)\).
\item
  Las trazas verticales son parábolas; las trazas horizontales son elipses (o círculos si \(a = b\)).
\item
  Es simétrica respecto a los planos \(xz\) e \(yz\), y respecto al eje \(z\).
\item
  \textbf{No es una superficie reglada}.
\end{itemize}

🎧 Podcast: 05\_cuadricas\_propiedades-geometricas\_3

Tu navegador no soporta el elemento de audio.

\subsection{Ejemplo resuelto}\label{ejemplo-resuelto-19}

\begin{quote}
\textbf{Ejemplo 4.} Identifique y describa la superficie \(z = 2x^2 + y^2\).
\end{quote}

\textbf{Solución.} La ecuación puede escribirse como:

\[
\frac{x^2}{(1/\sqrt{2})^2} + \frac{y^2}{1^2} = z
\]

Es un \textbf{paraboloide elíptico} que abre hacia arriba. Su vértice está en el origen. Las trazas horizontales \(z = k\) (con \(k > 0\)) son elipses. Las trazas verticales \(x = 0\) y \(y = 0\) son parábolas \(z = y^2\) y \(z = 2x^2\) respectivamente.

\label{fig:paraboloide-eliptico-ejemplo}Paraboloide elíptico del Ejemplo 4.

\begin{center}\rule{0.5\linewidth}{0.5pt}\end{center}

\section{Paraboloide hiperbólico}\label{paraboloide-hiperbuxf3lico}

\subsection{Ecuación canónica}\label{ecuaciuxf3n-canuxf3nica-4}

\[
\frac{x^2}{a^2} - \frac{y^2}{b^2} = \frac{z}{c}
\]

\subsection{Trazas}\label{trazas-4}

\begin{longtable}[]{@{}
  >{\raggedright\arraybackslash}p{(\columnwidth - 4\tabcolsep) * \real{0.2000}}
  >{\raggedright\arraybackslash}p{(\columnwidth - 4\tabcolsep) * \real{0.6286}}
  >{\raggedright\arraybackslash}p{(\columnwidth - 4\tabcolsep) * \real{0.1714}}@{}}
\toprule\noalign{}
\begin{minipage}[b]{\linewidth}\raggedright
Plano
\end{minipage} & \begin{minipage}[b]{\linewidth}\raggedright
Ecuación de la traza
\end{minipage} & \begin{minipage}[b]{\linewidth}\raggedright
Tipo
\end{minipage} \\
\midrule\noalign{}
\endhead
\bottomrule\noalign{}
\endlastfoot
\(xy\) (\(z = 0\)) & \(\frac{x^2}{a^2} - \frac{y^2}{b^2} = 0\) & Dos rectas que se cruzan en el origen \\
\(xz\) (\(y = 0\)) & \(z = \frac{c}{a^2} x^2\) & Parábola \\
\(yz\) (\(x = 0\)) & \(z = -\frac{c}{b^2} y^2\) & Parábola (abre en dirección opuesta) \\
\(z = k\) (\(k \neq 0\)) & \(\frac{x^2}{a^2} - \frac{y^2}{b^2} = \frac{k}{c}\) & Hipérbola \\
\end{longtable}

\subsection{Propiedades geométricas}\label{propiedades-geomuxe9tricas-5}

\begin{itemize}
\tightlist
\item
  Es una superficie \textbf{no acotada}.
\item
  También conocida como \textbf{silla de montar} por su forma característica.
\item
  Tiene un \textbf{punto de silla} en el origen.
\item
  Es una superficie \textbf{reglada}: contiene dos familias de rectas.
\item
  Es simétrica respecto al origen, a los planos \(xz\) e \(yz\), y al eje \(z\).
\end{itemize}

La propiedad de ser reglada es particularmente interesante: aunque la superficie tiene curvatura negativa (es una silla), puede construirse físicamente con rectas rígidas. Esto explica su uso en arquitectura para techos tensados.

🎧 Podcast: 05\_cuadricas\_propiedades-geometricas\_4

Tu navegador no soporta el elemento de audio.

\subsection{Ejemplo resuelto}\label{ejemplo-resuelto-20}

\begin{quote}
\textbf{Ejemplo 5.} Analice la superficie \(z = y^2 - x^2\).
\end{quote}

\textbf{Solución.} Se trata de un \textbf{paraboloide hiperbólico} con \(a = b = c = 1\). En el plano \(xy\) (\(z = 0\)), la ecuación se factoriza como \((y - x)(y + x) = 0\), que representa dos rectas \(y = x\) y \(y = -x\). Las trazas en planos \(x = 0\) y \(y = 0\) son parábolas \(z = y^2\) (abre hacia arriba) y \(z = -x^2\) (abre hacia abajo) respectivamente. Las trazas horizontales \(z = k\) son hipérbolas.

\label{fig:paraboloide-hiperbolico-ejemplo}Paraboloide hiperbólico del Ejemplo 5 (silla de montar).

\begin{center}\rule{0.5\linewidth}{0.5pt}\end{center}

\section{Cono elíptico}\label{cono-eluxedptico}

\subsection{Ecuación canónica}\label{ecuaciuxf3n-canuxf3nica-5}

\[
\frac{x^2}{a^2} + \frac{y^2}{b^2} = \frac{z^2}{c^2}
\]

\subsection{Trazas}\label{trazas-5}

\begin{longtable}[]{@{}
  >{\raggedright\arraybackslash}p{(\columnwidth - 4\tabcolsep) * \real{0.2000}}
  >{\raggedright\arraybackslash}p{(\columnwidth - 4\tabcolsep) * \real{0.6286}}
  >{\raggedright\arraybackslash}p{(\columnwidth - 4\tabcolsep) * \real{0.1714}}@{}}
\toprule\noalign{}
\begin{minipage}[b]{\linewidth}\raggedright
Plano
\end{minipage} & \begin{minipage}[b]{\linewidth}\raggedright
Ecuación de la traza
\end{minipage} & \begin{minipage}[b]{\linewidth}\raggedright
Tipo
\end{minipage} \\
\midrule\noalign{}
\endhead
\bottomrule\noalign{}
\endlastfoot
\(xy\) (\(z = 0\)) & \(\frac{x^2}{a^2} + \frac{y^2}{b^2} = 0\) & Punto \((0, 0, 0)\) \\
\(xz\) (\(y = 0\)) & \(z = \pm \frac{c}{a} x\) & Dos rectas que se cruzan \\
\(yz\) (\(x = 0\)) & \(z = \pm \frac{c}{b} y\) & Dos rectas que se cruzan \\
\(z = k\) (\(k \neq 0\)) & \(\frac{x^2}{a^2} + \frac{y^2}{b^2} = \frac{k^2}{c^2}\) & Elipse \\
\end{longtable}

\subsection{Propiedades geométricas}\label{propiedades-geomuxe9tricas-6}

\begin{itemize}
\tightlist
\item
  Es una superficie \textbf{no acotada}.
\item
  Tiene un \textbf{vértice} en el origen.
\item
  Es una superficie \textbf{reglada} (de hecho, está formada por rectas que pasan por el vértice).
\item
  Es simétrica respecto a todos los planos coordenados, ejes coordenados y el origen.
\item
  Cuando \(a = b\), se denomina \textbf{cono circular recto}.
\end{itemize}

🎧 Podcast: 05\_cuadricas\_propiedades-geometricas\_5

Tu navegador no soporta el elemento de audio.

\subsection{Ejemplo resuelto}\label{ejemplo-resuelto-21}

\begin{quote}
\textbf{Ejemplo 6.} Identifique la superficie \(4x^2 + y^2 - z^2 = 0\).
\end{quote}

\textbf{Solución.} Reescribiendo:

\[
\frac{x^2}{(1/2)^2} + \frac{y^2}{1^2} = \frac{z^2}{1^2}
\]

Es un \textbf{cono elíptico} con vértice en el origen. Las trazas horizontales \(z = k\) son elipses. Si cortamos el cono con un plano \(z = k\) con \(k \neq 0\), la elipse tiene semiejes \(a|k|/c = |k|/2\) y \(b|k|/c = |k|\).

\label{fig:cono-ejemplo}Cono elíptico del Ejemplo 6.

\begin{center}\rule{0.5\linewidth}{0.5pt}\end{center}

\section{Clasificación resumida de las cuádricas}\label{clasificaciuxf3n-resumida-de-las-cuuxe1dricas}

La siguiente tabla resume las seis superficies cuádricas principales en su forma canónica centrada en el origen:

\begin{longtable}[]{@{}
  >{\raggedright\arraybackslash}p{(\columnwidth - 8\tabcolsep) * \real{0.1739}}
  >{\raggedright\arraybackslash}p{(\columnwidth - 8\tabcolsep) * \real{0.2754}}
  >{\raggedright\arraybackslash}p{(\columnwidth - 8\tabcolsep) * \real{0.2899}}
  >{\raggedright\arraybackslash}p{(\columnwidth - 8\tabcolsep) * \real{0.1304}}
  >{\raggedright\arraybackslash}p{(\columnwidth - 8\tabcolsep) * \real{0.1304}}@{}}
\toprule\noalign{}
\begin{minipage}[b]{\linewidth}\raggedright
Superficie
\end{minipage} & \begin{minipage}[b]{\linewidth}\raggedright
Ecuación canónica
\end{minipage} & \begin{minipage}[b]{\linewidth}\raggedright
Signos cuadráticos
\end{minipage} & \begin{minipage}[b]{\linewidth}\raggedright
Acotada
\end{minipage} & \begin{minipage}[b]{\linewidth}\raggedright
Reglada
\end{minipage} \\
\midrule\noalign{}
\endhead
\bottomrule\noalign{}
\endlastfoot
Elipsoide & \(\frac{x^2}{a^2} + \frac{y^2}{b^2} + \frac{z^2}{c^2} = 1\) & \(+, +, +\) & Sí & No \\
Hiperboloide (1 hoja) & \(\frac{x^2}{a^2} + \frac{y^2}{b^2} - \frac{z^2}{c^2} = 1\) & \(+, +, -\) & No & Sí \\
Hiperboloide (2 hojas) & \(\frac{x^2}{a^2} + \frac{y^2}{b^2} - \frac{z^2}{c^2} = -1\) & \(+, +, -\) & No & No \\
Paraboloide elíptico & \(\frac{x^2}{a^2} + \frac{y^2}{b^2} = z\) & \(+, +\) (lineal) & No & No \\
Paraboloide hiperbólico & \(\frac{x^2}{a^2} - \frac{y^2}{b^2} = z\) & \(+, -\) (lineal) & No & Sí \\
Cono elíptico & \(\frac{x^2}{a^2} + \frac{y^2}{b^2} = z^2\) & \(+, +, -\) & No & Sí \\
\end{longtable}

\subsection{Criterio de clasificación}\label{criterio-de-clasificaciuxf3n-1}

Para clasificar una cuádrica dada por su ecuación canónica centrada en el origen:

\begin{enumerate}
\def\labelenumi{\arabic{enumi}.}
\tightlist
\item
  \textbf{Si todos los signos son positivos} \(\rightarrow\) Elipsoide (o esfera si \(a = b = c\)).
\item
  \textbf{Si hay dos signos positivos y uno negativo}:

  \begin{itemize}
  \tightlist
  \item
    Igual a \(+1\) \(\rightarrow\) Hiperboloide de una hoja.
  \item
    Igual a \(-1\) \(\rightarrow\) Hiperboloide de dos hojas.
  \end{itemize}
\item
  \textbf{Si hay una variable lineal} (sin cuadrado) \(\rightarrow\) Paraboloide.

  \begin{itemize}
  \tightlist
  \item
    Si los dos cuadrados tienen el mismo signo \(\rightarrow\) Paraboloide elíptico.
  \item
    Si los dos cuadrados tienen signos opuestos \(\rightarrow\) Paraboloide hiperbólico.
  \end{itemize}
\item
  \textbf{Si el lado derecho es cero} y hay dos cuadrados positivos y uno negativo \(\rightarrow\) Cono elíptico.
\end{enumerate}

\begin{center}\rule{0.5\linewidth}{0.5pt}\end{center}

\section{Ejercicios propuestos}\label{ejercicios-propuestos-5}

\begin{enumerate}
\def\labelenumi{\arabic{enumi}.}
\tightlist
\item
  Identifique y describa las trazas de la superficie \(9x^2 + 4y^2 + z^2 = 36\).
\item
  Determine qué tipo de cuádrica representa \(x^2 - y^2 + z^2 = 1\) y bosqueje sus trazas principales.
\item
  Halle la ecuación del elipsoide con vértices en \((\pm 3, 0, 0)\), \((0, \pm 2, 0)\) y \((0, 0, \pm 5)\).
\item
  Demuestre que el hiperboloide de una hoja \(x^2 + y^2 - z^2 = 1\) contiene las rectas parametrizadas por \(\mathbf{r}(t) = (\cos\theta - t\sin\theta, \sin\theta + t\cos\theta, t)\) para cualquier ángulo fijo \(\theta\).
\item
  Explique por qué el paraboloide hiperbólico \(z = x^2 - y^2\) es una superficie reglada, encontrando explícitamente dos familias de rectas contenidas en ella.
\end{enumerate}

\chapter{Integrales Múltiples}\label{integrales}

🎧 Podcast del capítulo

Tu navegador no soporta el elemento de audio.

El cálculo integral se extiende de manera natural a funciones de varias variables mediante las integrales dobles y triples. Estas herramientas permiten calcular volúmenes, áreas, centroides, momentos de inercia y promedios sobre regiones bidimensionales y tridimensionales.

\section{Integrales dobles}\label{integrales-dobles}

\subsection{Definición por sumas de Riemann}\label{definiciuxf3n-por-sumas-de-riemann}

Sea \(f(x, y)\) una función definida sobre una región rectangular \(R = [a, b] \times [c, d]\). Dividimos \(R\) en \(m \times n\) subrectángulos de área \(\Delta A = \Delta x \Delta y\). La \textbf{integral doble} de \(f\) sobre \(R\) es:

\[
\iint_R f(x, y) \, dA = \lim_{m,n \to \infty} \sum_{i=1}^m \sum_{j=1}^n f(x_i^*, y_j^*) \Delta A
\]

siempre que el límite exista.

🎧 Podcast: 06\_integrales\_definicion-por-sumas-de-riemann

Tu navegador no soporta el elemento de audio.

\subsection{Propiedades}\label{propiedades-1}

\begin{enumerate}
\def\labelenumi{\arabic{enumi}.}
\tightlist
\item
  \textbf{Linealidad}: \(\iint_R [\alpha f + \beta g] \, dA = \alpha \iint_R f \, dA + \beta \iint_R g \, dA\)
\item
  \textbf{Monotonía}: Si \(f(x, y) \leq g(x, y)\), entonces \(\iint_R f \, dA \leq \iint_R g \, dA\)
\item
  \textbf{Aditividad}: Si \(R = R_1 \cup R_2\) con interiores disjuntos, \(\iint_R f \, dA = \iint_{R_1} f \, dA + \iint_{R_2} f \, dA\)
\item
  \textbf{Área}: \(\iint_R 1 \, dA = A(R)\)
\end{enumerate}

🎧 Podcast: 06\_integrales\_propiedades

Tu navegador no soporta el elemento de audio.

\subsection{Evaluación: integrales iteradas (Teorema de Fubini)}\label{evaluaciuxf3n-integrales-iteradas-teorema-de-fubini}

Si \(f\) es continua sobre \(R = [a, b] \times [c, d]\), entonces:

\[
\iint_R f(x, y) \, dA = \int_a^b \int_c^d f(x, y) \, dy \, dx = \int_c^d \int_a^b f(x, y) \, dx \, dy
\]

El orden de integración puede intercambiarse cuando \(f\) es continua.

🎧 Podcast: 06\_integrales\_evaluacion-integrales-iteradas-teorema-de-fubini

Tu navegador no soporta el elemento de audio.

\subsection{Regiones generales}\label{regiones-generales}

\subsubsection{Tipo I (verticalmente simple)}\label{tipo-i-verticalmente-simple}

\[
R = \{(x, y) : a \leq x \leq b, \; g_1(x) \leq y \leq g_2(x)\}
\]

\[
\iint_R f(x, y) \, dA = \int_a^b \int_{g_1(x)}^{g_2(x)} f(x, y) \, dy \, dx
\]

\subsubsection{Tipo II (horizontalmente simple)}\label{tipo-ii-horizontalmente-simple}

\[
R = \{(x, y) : c \leq y \leq d, \; h_1(y) \leq x \leq h_2(y)\}
\]

\[
\iint_R f(x, y) \, dA = \int_c^d \int_{h_1(y)}^{h_2(y)} f(x, y) \, dx \, dy
\]

\subsection{Ejemplo resuelto}\label{ejemplo-resuelto-22}

\begin{quote}
\textbf{Ejemplo 1.} Evalúe \(\iint_R (x + y) \, dA\) donde \(R\) está acotada por \(y = x^2\) e \(y = x\).
\end{quote}

\textbf{Solución.} Las curvas se intersectan cuando \(x^2 = x\), es decir, \(x = 0\) y \(x = 1\).

Región Tipo I: \(0 \leq x \leq 1\), \(x^2 \leq y \leq x\).

\[
\begin{aligned}
\iint_R (x + y) \, dA &= \int_0^1 \int_{x^2}^x (x + y) \, dy \, dx \\
&= \int_0^1 \left[ xy + \frac{y^2}{2} \right]_{x^2}^x dx \\
&= \int_0^1 \left( x^2 + \frac{x^2}{2} - x^3 - \frac{x^4}{2} \right) dx \\
&= \int_0^1 \left( \frac{3x^2}{2} - x^3 - \frac{x^4}{2} \right) dx \\
&= \left[ \frac{x^3}{2} - \frac{x^4}{4} - \frac{x^5}{10} \right]_0^1 = \frac{1}{2} - \frac{1}{4} - \frac{1}{10} = \frac{3}{20}
\end{aligned}
\]

\section{Cambio de variables en integrales dobles}\label{cambio-de-variables-en-integrales-dobles}

\subsection{Transformación general}\label{transformaciuxf3n-general}

Si \(x = x(u, v)\), \(y = y(u, v)\) es una transformación biyectiva con jacobiano:

\[
J = \frac{\partial(x, y)}{\partial(u, v)} = \begin{vmatrix}
\frac{\partial x}{\partial u} & \frac{\partial x}{\partial v} \\
\frac{\partial y}{\partial u} & \frac{\partial y}{\partial v}
\end{vmatrix}
\]

Entonces:

\[
\iint_R f(x, y) \, dA = \iint_S f(x(u, v), y(u, v)) \, |J| \, du \, dv
\]

\subsection{Coordenadas polares}\label{coordenadas-polares}

En coordenadas polares: \(x = r\cos\theta\), \(y = r\sin\theta\).

\[
J = \begin{vmatrix}
\cos\theta & -r\sin\theta \\
\sin\theta & r\cos\theta
\end{vmatrix} = r
\]

Por tanto:

\[
\iint_R f(x, y) \, dA = \iint_S f(r\cos\theta, r\sin\theta) \, r \, dr \, d\theta
\]

\subsection{Ejemplo resuelto}\label{ejemplo-resuelto-23}

\begin{quote}
\textbf{Ejemplo 2.} Evalúe \(\iint_D e^{x^2 + y^2} \, dA\) donde \(D\) es el disco unitario.
\end{quote}

\textbf{Solución.} En coordenadas polares: \(0 \leq r \leq 1\), \(0 \leq \theta \leq 2\pi\).

\[
\begin{aligned}
\iint_D e^{x^2 + y^2} \, dA &= \int_0^{2\pi} \int_0^1 e^{r^2} r \, dr \, d\theta \\
&= 2\pi \cdot \frac{1}{2} e^{r^2} \Big|_0^1 = \pi(e - 1)
\end{aligned}
\]

\section{Integrales triples}\label{integrales-triples}

\subsection{Definición}\label{definiciuxf3n-5}

La \textbf{integral triple} de \(f(x, y, z)\) sobre una región sólida \(E\) es:

\[
\iiint_E f(x, y, z) \, dV = \lim_{l,m,n \to \infty} \sum_{i=1}^l \sum_{j=1}^m \sum_{k=1}^n f(x_i^*, y_j^*, z_k^*) \Delta V
\]

🎧 Podcast: 06\_integrales\_definicion

Tu navegador no soporta el elemento de audio.

\subsection{Evaluación como integral iterada}\label{evaluaciuxf3n-como-integral-iterada}

Para una región \(E = \{(x, y, z) : a \leq x \leq b, \; g_1(x) \leq y \leq g_2(x), \; u_1(x, y) \leq z \leq u_2(x, y)\}\):

\[
\iiint_E f(x, y, z) \, dV = \int_a^b \int_{g_1(x)}^{g_2(x)} \int_{u_1(x,y)}^{u_2(x,y)} f(x, y, z) \, dz \, dy \, dx
\]

Existen 6 órdenes posibles de integración.

\subsection{Ejemplo resuelto}\label{ejemplo-resuelto-24}

\begin{quote}
\textbf{Ejemplo 3.} Calcule el volumen del tetraedro acotado por los planos coordenados y \(x + y + z = 1\).
\end{quote}

\textbf{Solución.}

\[
\begin{aligned}
V &= \int_0^1 \int_0^{1-x} \int_0^{1-x-y} dz \, dy \, dx \\
&= \int_0^1 \int_0^{1-x} (1 - x - y) \, dy \, dx \\
&= \int_0^1 \left[ y - xy - \frac{y^2}{2} \right]_0^{1-x} dx \\
&= \int_0^1 \left( (1-x) - x(1-x) - \frac{(1-x)^2}{2} \right) dx \\
&= \int_0^1 \frac{(1-x)^2}{2} dx = \frac{1}{2} \left[ -\frac{(1-x)^3}{3} \right]_0^1 = \frac{1}{6}
\end{aligned}
\]

\section{Cambio de variables en integrales triples}\label{cambio-de-variables-en-integrales-triples}

\subsection{Coordenadas cilíndricas}\label{coordenadas-ciluxedndricas-1}

\[
\iiint_E f(x, y, z) \, dV = \iiint_E f(r\cos\theta, r\sin\theta, z) \, r \, dr \, d\theta \, dz
\]

\subsection{Coordenadas esféricas}\label{coordenadas-esfuxe9ricas-1}

\[
\iiint_E f(x, y, z) \, dV = \iiint_E f(\rho\sin\phi\cos\theta, \rho\sin\phi\sin\theta, \rho\cos\phi) \, \rho^2 \sin\phi \, d\rho \, d\phi \, d\theta
\]

\subsection{Ejemplo resuelto}\label{ejemplo-resuelto-25}

\begin{quote}
\textbf{Ejemplo 4.} Calcule el volumen de una esfera de radio \(R\) usando coordenadas esféricas.
\end{quote}

\textbf{Solución.}

\[
\begin{aligned}
V &= \int_0^{2\pi} \int_0^{\pi} \int_0^R \rho^2 \sin\phi \, d\rho \, d\phi \, d\theta \\
&= \int_0^{2\pi} d\theta \int_0^{\pi} \sin\phi \, d\phi \int_0^R \rho^2 \, d\rho \\
&= (2\pi)(2)\left(\frac{R^3}{3}\right) = \frac{4}{3}\pi R^3
\end{aligned}
\]

\section{Aplicaciones de las integrales múltiples}\label{aplicaciones-de-las-integrales-muxfaltiples}

\subsection{Áreas y volúmenes}\label{uxe1reas-y-voluxfamenes}

\begin{itemize}
\tightlist
\item
  Área de una región plana: \(A = \iint_R 1 \, dA\)
\item
  Volumen bajo una superficie: \(V = \iint_R f(x, y) \, dA\) (con \(f \geq 0\))
\item
  Volumen de un sólido: \(V = \iiint_E 1 \, dV\)
\end{itemize}

\subsection{Masa y centro de masa}\label{masa-y-centro-de-masa}

Para una lámina con densidad \(\rho(x, y)\):

\[
M = \iint_R \rho(x, y) \, dA
\]

\[
\bar{x} = \frac{1}{M} \iint_R x \rho(x, y) \, dA, \qquad
\bar{y} = \frac{1}{M} \iint_R y \rho(x, y) \, dA
\]

\subsection{Momentos de inercia}\label{momentos-de-inercia}

\[
I_x = \iint_R y^2 \rho(x, y) \, dA, \qquad
I_y = \iint_R x^2 \rho(x, y) \, dA, \qquad
I_0 = I_x + I_y
\]

\subsection{Valor promedio}\label{valor-promedio}

El valor promedio de \(f\) sobre \(R\) es:

\[
f_{\text{prom}} = \frac{1}{A(R)} \iint_R f(x, y) \, dA
\]

\section{Visualización interactiva: regiones de integración}\label{visualizaciuxf3n-interactiva-regiones-de-integraciuxf3n}

Explore el siguiente applet de GeoGebra para visualizar regiones de integración en coordenadas cartesianas y polares:

\section{Ejercicios propuestos}\label{ejercicios-propuestos-6}

\begin{enumerate}
\def\labelenumi{\arabic{enumi}.}
\tightlist
\item
  Evalúe \(\int_0^1 \int_0^{x^2} (x + y) \, dy \, dx\).
\item
  Cambie el orden de integración y evalúe \(\int_0^1 \int_y^1 e^{x^2} \, dx \, dy\).
\item
  Use coordenadas polares para calcular \(\iint_D \sqrt{x^2 + y^2} \, dA\) donde \(D\) es el disco \(x^2 + y^2 \leq 4\).
\item
  Calcule el volumen del sólido acotado por el paraboloide \(z = x^2 + y^2\) y el plano \(z = 4\).
\item
  Halle el centro de masa de una lámina semicircular de radio \(R\) con densidad uniforme.
\item
  Use coordenadas esféricas para calcular \(\iiint_E (x^2 + y^2 + z^2) \, dV\) donde \(E\) es la bola unitaria.
\end{enumerate}

\chapter{Geometría Diferencial de Superficies}\label{diferencial}

🎧 Podcast del capítulo

Tu navegador no soporta el elemento de audio.

En este capítulo se desarrollan las herramientas del cálculo diferencial para el estudio de las superficies. Estas técnicas permiten analizar propiedades locales como la orientación, el plano tangente, la curvatura y la métrica de la superficie.

\section{Superficies paramétricas}\label{superficies-paramuxe9tricas}

\subsection{Definición}\label{definiciuxf3n-6}

Una \textbf{superficie paramétrica} es una aplicación diferenciable:

\[
\mathbf{r}: D \subseteq \mathbb{R}^2 \longrightarrow \mathbb{R}^3, \qquad
\mathbf{r}(u, v) = \bigl( x(u, v),\, y(u, v),\, z(u, v) \bigr)
\]

Se requiere que los vectores tangentes parciales:

\[
\mathbf{r}_u = \frac{\partial \mathbf{r}}{\partial u}, \qquad
\mathbf{r}_v = \frac{\partial \mathbf{r}}{\partial v}
\]

sean linealmente independientes en cada punto, es decir, \(\mathbf{r}_u \times \mathbf{r}_v \neq \mathbf{0}\). Esta condición garantiza que la superficie tenga un plano tangente bien definido en cada punto.

🎧 Podcast: 07\_diferencial\_definicion

Tu navegador no soporta el elemento de audio.

\subsection{Ejemplos de parametrizaciones}\label{ejemplos-de-parametrizaciones}

\textbf{Esfera de radio \(R\):}

\[
\mathbf{r}(\phi, \theta) = \bigl( R\sin\phi\cos\theta,\; R\sin\phi\sin\theta,\; R\cos\phi \bigr)
\]

con \(0 \leq \phi \leq \pi\) y \(0 \leq \theta < 2\pi\).

\textbf{Elipsoide:}

\[
\mathbf{r}(\phi, \theta) = \bigl( a\sin\phi\cos\theta,\; b\sin\phi\sin\theta,\; c\cos\phi \bigr)
\]

\textbf{Cilindro circular de radio \(r\):}

\[
\mathbf{r}(\theta, z) = \bigl( r\cos\theta,\; r\sin\theta,\; z \bigr)
\]

\textbf{Paraboloide hiperbólico:}

\[
\mathbf{r}(u, v) = \left( u,\; v,\; \frac{u^2}{a^2} - \frac{v^2}{b^2} \right)
\]

\textbf{Hiperboloide de una hoja:}

\[
\mathbf{r}(u, v) = \bigl( a\cosh u \cos v,\; b\cosh u \sin v,\; c\sinh u \bigr)
\]

\section{Cambio de parametrización}\label{cambio-de-parametrizaciuxf3n}

Dos parametrizaciones \(\mathbf{r}(u, v)\) y \(\mathbf{\tilde{r}}(\tilde{u}, \tilde{v})\) describen el mismo trozo de superficie si existe un difeomorfismo (cambio de coordenadas suave e invertible) \((u, v) = \Phi(\tilde{u}, \tilde{v})\) tal que:

\[
\mathbf{\tilde{r}}(\tilde{u}, \tilde{v}) = \mathbf{r}\bigl( \Phi(\tilde{u}, \tilde{v}) \bigr)
\]

Las propiedades geométricas intrínsecas de la superficie (longitudes, ángulos, áreas, curvatura gaussiana) no dependen de la parametrización elegida.

\section{Vector normal y orientación}\label{vector-normal-y-orientaciuxf3n}

\subsection{Vector normal}\label{vector-normal}

El \textbf{vector normal} a la superficie en un punto está dado por el producto cruz de los vectores tangentes:

\[
\mathbf{N}(u, v) = \mathbf{r}_u \times \mathbf{r}_v
\]

El vector unitario normal es:

\[
\mathbf{n}(u, v) = \frac{\mathbf{r}_u \times \mathbf{r}_v}{\|\mathbf{r}_u \times \mathbf{r}_v\|}
\]

\subsection{Interpretación geométrica}\label{interpretaciuxf3n-geomuxe9trica-2}

El vector normal es perpendicular al plano tangente y determina la \textbf{orientación} de la superficie. Una superficie es \textbf{orientable} si existe un campo continuo de vectores normales unitarios definido globalmente.

\begin{quote}
\textbf{Ejemplo.} La esfera es orientable. El plano tangente en cada punto es perpendicular al radio vector. Una superficie como la cinta de Möbius no es orientable.
\end{quote}

\subsection{Ejemplo de cálculo}\label{ejemplo-de-cuxe1lculo}

\begin{quote}
\textbf{Ejemplo.} Calcule el vector normal unitario al paraboloide elíptico \(\mathbf{r}(u, v) = (u, v, u^2 + v^2)\).
\end{quote}

\textbf{Solución.} Los vectores tangentes son:

\[
\mathbf{r}_u = (1, 0, 2u), \qquad \mathbf{r}_v = (0, 1, 2v)
\]

El producto cruz es:

\[
\mathbf{r}_u \times \mathbf{r}_v = \begin{vmatrix}
\mathbf{i} & \mathbf{j} & \mathbf{k} \\
1 & 0 & 2u \\
0 & 1 & 2v
\end{vmatrix} = (-2u, -2v, 1)
\]

Su magnitud es \(\|\mathbf{r}_u \times \mathbf{r}_v\| = \sqrt{4u^2 + 4v^2 + 1}\). Por tanto:

\[
\mathbf{n} = \frac{(-2u, -2v, 1)}{\sqrt{4u^2 + 4v^2 + 1}}
\]

\section{Plano tangente}\label{plano-tangente}

\subsection{Definición}\label{definiciuxf3n-7}

El \textbf{plano tangente} a la superficie en el punto \(\mathbf{r}(u_0, v_0)\) es el plano que pasa por dicho punto y es perpendicular al vector normal \(\mathbf{N}(u_0, v_0)\).

🎧 Podcast: 07\_diferencial\_definicion\_1

Tu navegador no soporta el elemento de audio.

\subsection{Ecuación del plano tangente}\label{ecuaciuxf3n-del-plano-tangente-1}

Si \(\mathbf{x} = (x, y, z)\) es un punto genérico del plano tangente, entonces:

\[
\bigl( \mathbf{r}_u \times \mathbf{r}_v \bigr) \cdot \bigl( \mathbf{x} - \mathbf{r}(u_0, v_0) \bigr) = 0
\]

Equivalentemente, usando el vector unitario normal:

\[
\mathbf{n} \cdot \bigl( \mathbf{x} - \mathbf{r}(u_0, v_0) \bigr) = 0
\]

\subsection{Ejemplo}\label{ejemplo-2}

\begin{quote}
\textbf{Ejemplo.} Halle la ecuación del plano tangente al elipsoide \(\frac{x^2}{4} + \frac{y^2}{9} + z^2 = 1\) en el punto \((1, 2, \frac{\sqrt{2}}{2})\).
\end{quote}

\textbf{Solución.} Sea \(F(x, y, z) = \frac{x^2}{4} + \frac{y^2}{9} + z^2 - 1\). El gradiente \(\nabla F\) es normal a la superficie:

\[
\nabla F = \left( \frac{x}{2}, \frac{2y}{9}, 2z \right)
\]

En el punto dado:

\[
\nabla F(1, 2, \tfrac{\sqrt{2}}{2}) = \left( \frac{1}{2}, \frac{4}{9}, \sqrt{2} \right)
\]

La ecuación del plano tangente es:

\[
\frac{1}{2}(x - 1) + \frac{4}{9}(y - 2) + \sqrt{2}\left(z - \frac{\sqrt{2}}{2}\right) = 0
\]

Simplificando:

\[
\frac{x}{2} + \frac{4y}{9} + \sqrt{2}\,z = 1 + \frac{8}{9} + 1 = \frac{26}{9}
\]

Multiplicando por 18:

\[
9x + 8y + 18\sqrt{2}\,z = 52
\]

\section{Primera forma fundamental}\label{primera-forma-fundamental}

La primera forma fundamental codifica la métrica intrínseca de la superficie: permite calcular longitudes de arco, ángulos y áreas sobre la superficie usando solo las coordenadas \((u, v)\).

\subsection{Coeficientes fundamentales}\label{coeficientes-fundamentales}

\[
\begin{aligned}
E &= \mathbf{r}_u \cdot \mathbf{r}_u = \|\mathbf{r}_u\|^2 \\
F &= \mathbf{r}_u \cdot \mathbf{r}_v \\
G &= \mathbf{r}_v \cdot \mathbf{r}_v = \|\mathbf{r}_v\|^2
\end{aligned}
\]

Estos tres coeficientes constituyen la \textbf{primera forma fundamental} \(I = E\,du^2 + 2F\,du\,dv + G\,dv^2\).

\subsection{Aplicaciones}\label{aplicaciones}

\textbf{Elemento de longitud de arco:}

\[
ds^2 = E\,du^2 + 2F\,du\,dv + G\,dv^2
\]

\textbf{Área de una región \(R \subseteq D\):}

\[
A(R) = \iint_R \sqrt{EG - F^2} \, du\, dv = \iint_R \|\mathbf{r}_u \times \mathbf{r}_v\| \, du\, dv
\]

\textbf{Ángulo entre dos curvas} sobre la superficie: si las curvas se intersectan en un punto con vectores tangentes \((du_1, dv_1)\) y \((du_2, dv_2)\) en el espacio de parámetros, el ángulo \(\alpha\) entre ellas satisface:

\[
\cos\alpha = \frac{E\,du_1\,du_2 + F(du_1\,dv_2 + du_2\,dv_1) + G\,dv_1\,dv_2}{\sqrt{E\,du_1^2 + 2F\,du_1\,dv_1 + G\,dv_1^2} \; \sqrt{E\,du_2^2 + 2F\,du_2\,dv_2 + G\,dv_2^2}}
\]

\subsection{Ejemplo: Esfera}\label{ejemplo-esfera}

Para la esfera \(\mathbf{r}(\phi, \theta) = (R\sin\phi\cos\theta, R\sin\phi\sin\theta, R\cos\phi)\):

\[
\mathbf{r}_\phi = (R\cos\phi\cos\theta, R\cos\phi\sin\theta, -R\sin\phi), \quad
\mathbf{r}_\theta = (-R\sin\phi\sin\theta, R\sin\phi\cos\theta, 0)
\]

Entonces:

\[
\begin{aligned}
E &= \mathbf{r}_\phi \cdot \mathbf{r}_\phi = R^2 \\
F &= \mathbf{r}_\phi \cdot \mathbf{r}_\theta = 0 \\
G &= \mathbf{r}_\theta \cdot \mathbf{r}_\theta = R^2\sin^2\phi
\end{aligned}
\]

El elemento de área es \(\sqrt{EG - F^2} = R^2\sin\phi\), de donde el área total es:

\[
A = \int_0^{2\pi} \int_0^{\pi} R^2\sin\phi \, d\phi\, d\theta = 4\pi R^2
\]

\section{Segunda forma fundamental y curvatura}\label{segunda-forma-fundamental-y-curvatura}

La segunda forma fundamental mide cómo se curva la superficie en el espacio ambiente. Se define a partir de la derivada del vector normal.

\subsection{Vector normal y derivadas}\label{vector-normal-y-derivadas}

Si \(\mathbf{n}\) es el vector unitario normal, definimos:

\[
\begin{aligned}
L &= \mathbf{r}_{uu} \cdot \mathbf{n} = -\mathbf{r}_u \cdot \mathbf{n}_u \\
M &= \mathbf{r}_{uv} \cdot \mathbf{n} = -\mathbf{r}_u \cdot \mathbf{n}_v = -\mathbf{r}_v \cdot \mathbf{n}_u \\
N &= \mathbf{r}_{vv} \cdot \mathbf{n} = -\mathbf{r}_v \cdot \mathbf{n}_v
\end{aligned}
\]

Estos coeficientes constituyen la \textbf{segunda forma fundamental} \(II = L\,du^2 + 2M\,du\,dv + N\,dv^2\).

\subsection{Curvaturas principales}\label{curvaturas-principales}

En cada punto de la superficie, la curvatura varía según la dirección. Las \textbf{curvaturas principales} \(\kappa_1\) y \(\kappa_2\) son el máximo y el mínimo valor de la curvatura normal en el punto. Se obtienen como los valores propios de la matriz:

\[
\begin{pmatrix}
E & F \\
F & G
\end{pmatrix}^{-1}
\begin{pmatrix}
L & M \\
M & N
\end{pmatrix}
\]

\subsection{Curvatura gaussiana}\label{curvatura-gaussiana}

La \textbf{curvatura gaussiana} \(K\) es el producto de las curvaturas principales:

\[
K = \kappa_1 \kappa_2 = \frac{LN - M^2}{EG - F^2}
\]

Su signo tiene un significado geométrico profundo:

\begin{itemize}
\tightlist
\item
  \(K > 0\): la superficie es localmente como una esfera o un elipsoide (curvatura positiva).
\item
  \(K < 0\): la superficie es localmente como una silla de montar (curvatura negativa).
\item
  \(K = 0\): la superficie es localmente cilíndrica o plana (curvatura nula).
\end{itemize}

\subsection{Curvatura media}\label{curvatura-media}

La \textbf{curvatura media} \(H\) es el promedio aritmético de las curvaturas principales:

\[
H = \frac{\kappa_1 + \kappa_2}{2} = \frac{EN + GL - 2FM}{2(EG - F^2)}
\]

Una superficie con \(H = 0\) en todos sus puntos se denomina \textbf{superficie mínima}. Ejemplos clásicos son el catenoide y la helicoide.

\subsection{Ejemplo: Paraboloide elíptico}\label{ejemplo-paraboloide-eluxedptico}

\begin{quote}
\textbf{Ejemplo.} Calcule la curvatura gaussiana y la curvatura media del paraboloide elíptico \(z = x^2 + y^2\) en el origen.
\end{quote}

\textbf{Solución.} Parametrizamos como \(\mathbf{r}(u, v) = (u, v, u^2 + v^2)\). Ya calculamos:

\[
\mathbf{r}_u = (1, 0, 2u), \quad \mathbf{r}_v = (0, 1, 2v), \quad \mathbf{n} = \frac{(-2u, -2v, 1)}{\sqrt{4u^2 + 4v^2 + 1}}
\]

En el origen \((u, v) = (0, 0)\):

\[
\mathbf{r}_u = (1, 0, 0), \quad \mathbf{r}_v = (0, 1, 0), \quad \mathbf{n} = (0, 0, 1)
\]

Las segundas derivadas son:

\[
\mathbf{r}_{uu} = (0, 0, 2), \quad \mathbf{r}_{uv} = (0, 0, 0), \quad \mathbf{r}_{vv} = (0, 0, 2)
\]

Por tanto:

\[
\begin{aligned}
E &= 1, & F &= 0, & G &= 1 \\
L &= 2, & M &= 0, & N &= 2
\end{aligned}
\]

La curvatura gaussiana es:

\[
K = \frac{LN - M^2}{EG - F^2} = \frac{4 - 0}{1 - 0} = 4
\]

La curvatura media es:

\[
H = \frac{EN + GL - 2FM}{2(EG - F^2)} = \frac{2 + 2}{2} = 2
\]

Como \(K > 0\), la superficie es localmente convexa en el origen (como una copa hacia arriba).

\subsection{Ejemplo: Paraboloide hiperbólico}\label{ejemplo-paraboloide-hiperbuxf3lico}

\begin{quote}
\textbf{Ejemplo.} Calcule la curvatura gaussiana del paraboloide hiperbólico \(z = x^2 - y^2\) en el origen.
\end{quote}

\textbf{Solución.} Con \(\mathbf{r}(u, v) = (u, v, u^2 - v^2)\), en el origen:

\[
\mathbf{r}_u = (1, 0, 0), \quad \mathbf{r}_v = (0, 1, 0), \quad \mathbf{n} = (0, 0, 1)
\]

\[
\mathbf{r}_{uu} = (0, 0, 2), \quad \mathbf{r}_{uv} = (0, 0, 0), \quad \mathbf{r}_{vv} = (0, 0, -2)
\]

Entonces \(E = G = 1\), \(F = 0\), \(L = 2\), \(M = 0\), \(N = -2\).

\[
K = \frac{(2)(-2) - 0}{1} = -4
\]

Como \(K < 0\), el origen es un \textbf{punto de silla}, consistente con la forma de ``silla de montar'' de esta superficie.

\section{Líneas de curvatura y direcciones principales}\label{luxedneas-de-curvatura-y-direcciones-principales}

Las \textbf{direcciones principales} son las direcciones en las cuales se alcanzan las curvaturas principales. Una \textbf{línea de curvatura} es una curva sobre la superficie cuya tangente en cada punto tiene una dirección principal.

En los puntos umbilicales, donde \(\kappa_1 = \kappa_2\), todas las direcciones son principales. Por ejemplo, todos los puntos de una esfera son umbilicales.

\section{Ecuaciones de Weingarten}\label{ecuaciones-de-weingarten}

Las derivadas del vector normal unitario están relacionadas con las derivadas de la parametrización mediante las \textbf{ecuaciones de Weingarten}:

\[
\begin{aligned}
\mathbf{n}_u &= \frac{FM - GL}{EG - F^2} \mathbf{r}_u + \frac{FL - EM}{EG - F^2} \mathbf{r}_v \\
\mathbf{n}_v &= \frac{FN - GM}{EG - F^2} \mathbf{r}_u + \frac{FM - EN}{EG - F^2} \mathbf{r}_v
\end{aligned}
\]

Estas ecuaciones muestran que la variación del vector normal está completamente determinada por la primera y segunda formas fundamentales.

\chapter{Aplicaciones de las Superficies}\label{aplicaciones}

🎧 Podcast del capítulo

Tu navegador no soporta el elemento de audio.

Las superficies en tres dimensiones tienen aplicaciones fundamentales en múltiples disciplinas científicas y de ingeniería. En este capítulo se presentan las aplicaciones más importantes, con énfasis en el cálculo integral sobre superficies y sus usos prácticos.

\section{Áreas de superficies}\label{uxe1reas-de-superficies}

\subsection{Fórmula general}\label{fuxf3rmula-general}

Para una superficie paramétrica \(\mathbf{r}(u, v)\) con \((u, v) \in D\), el área está dada por:

\[
A = \iint_D \|\mathbf{r}_u \times \mathbf{r}_v\| \, du\, dv = \iint_D \sqrt{EG - F^2} \, du\, dv
\]

Si la superficie está dada explícitamente como \(z = f(x, y)\) sobre una región \(R\) del plano \(xy\), entonces:

\[
A = \iint_R \sqrt{1 + \left(\frac{\partial f}{\partial x}\right)^2 + \left(\frac{\partial f}{\partial y}\right)^2} \, dx\, dy
\]

\subsection{Ejemplo resuelto: Área de un paraboloide}\label{ejemplo-resuelto-uxe1rea-de-un-paraboloide}

\begin{quote}
\textbf{Ejemplo 1.} Calcule el área de la porción del paraboloide \(z = x^2 + y^2\) que queda debajo del plano \(z = 4\).
\end{quote}

\textbf{Solución.} La superficie está definida sobre el disco \(R = \{(x, y) : x^2 + y^2 \leq 4\}\). Las derivadas parciales son:

\[
\frac{\partial z}{\partial x} = 2x, \qquad \frac{\partial z}{\partial y} = 2y
\]

El elemento de área es:

\[
\sqrt{1 + 4x^2 + 4y^2} = \sqrt{1 + 4(x^2 + y^2)}
\]

Usando coordenadas polares (\(x = r\cos\theta\), \(y = r\sin\theta\), \(dx\,dy = r\,dr\,d\theta\)):

\[
A = \int_0^{2\pi} \int_0^2 \sqrt{1 + 4r^2} \, r\,dr\,d\theta
\]

Haciendo \(u = 1 + 4r^2\), \(du = 8r\,dr\):

\[
A = 2\pi \int_1^{17} \sqrt{u} \, \frac{du}{8} = \frac{\pi}{4} \cdot \frac{2}{3} u^{3/2} \Big|_1^{17} = \frac{\pi}{6} \bigl(17\sqrt{17} - 1\bigr)
\]

Por tanto:

\[
A = \frac{\pi}{6}(17\sqrt{17} - 1) \approx 36.18 \text{ unidades}^2
\]

\section{Integrales de superficie de campos escalares}\label{integrales-de-superficie-de-campos-escalares}

Dada una función escalar \(f(x, y, z)\) definida sobre una superficie \(S\), la integral de superficie de \(f\) sobre \(S\) es:

\[
\iint_S f(x, y, z) \, dS = \iint_D f\bigl(\mathbf{r}(u, v)\bigr) \|\mathbf{r}_u \times \mathbf{r}_v\| \, du\, dv
\]

\subsection{Interpretaciones físicas}\label{interpretaciones-fuxedsicas}

\begin{itemize}
\tightlist
\item
  Si \(f\) representa la densidad superficial de masa, la integral da la \textbf{masa total} de una lámina con la forma de \(S\).
\item
  Si \(f\) representa la densidad de carga eléctrica, la integral da la \textbf{carga total} sobre \(S\).
\end{itemize}

\subsection{Ejemplo resuelto: Masa de una lámina esférica}\label{ejemplo-resuelto-masa-de-una-luxe1mina-esfuxe9rica}

\begin{quote}
\textbf{Ejemplo 2.} Una lámina esférica de radio \(R\) tiene densidad de masa proporcional a la distancia al plano ecuatorial, es decir, \(\rho(x, y, z) = k|z|\). Calcule la masa total de la lámina.
\end{quote}

\textbf{Solución.} Parametrizamos la esfera completa con coordenadas esféricas. El elemento de área es \(dS = R^2\sin\phi \, d\phi\, d\theta\) y \(z = R\cos\phi\). Entonces:

\[
M = \iint_S k|z| \, dS = kR \int_0^{2\pi} \int_0^{\pi} |R\cos\phi| \, R\sin\phi \, d\phi\, d\theta
\]

Separando en \(\phi \in [0, \pi/2]\) (donde \(\cos\phi \geq 0\)) y \(\phi \in [\pi/2, \pi]\) (donde \(\cos\phi \leq 0\)):

\[
M = 2\pi k R^3 \left[ \int_0^{\pi/2} \cos\phi\sin\phi \, d\phi + \int_{\pi/2}^{\pi} (-\cos\phi)\sin\phi \, d\phi \right]
\]

Cada integral vale \(\frac{1}{2}\), por tanto:

\[
M = 2\pi k R^3 \left( \frac{1}{2} + \frac{1}{2} \right) = 2\pi k R^3
\]

\section{Integrales de superficie de campos vectoriales (Flujo)}\label{integrales-de-superficie-de-campos-vectoriales-flujo}

Dado un campo vectorial \(\mathbf{F}(x, y, z)\) y una superficie orientable \(S\) con vector normal unitario \(\mathbf{n}\), el \textbf{flujo} de \(\mathbf{F}\) a través de \(S\) es:

\[
\Phi = \iint_S \mathbf{F} \cdot \mathbf{n} \, dS = \iint_S \mathbf{F} \cdot d\mathbf{S}
\]

En términos de una parametrización:

\[
\Phi = \iint_D \mathbf{F}\bigl(\mathbf{r}(u, v)\bigr) \cdot (\mathbf{r}_u \times \mathbf{r}_v) \, du\, dv
\]

\subsection{Interpretación física}\label{interpretaciuxf3n-fuxedsica}

El flujo representa la cantidad de ``fluido'' (material o abstracto) que atraviesa la superficie por unidad de tiempo. Si \(\mathbf{F}\) es un campo de velocidades de un fluido, el flujo es el caudal volumétrico a través de \(S\).

\subsection{Teorema de la divergencia (Gauss)}\label{teorema-de-la-divergencia-gauss}

Si \(V\) es un sólido acotado cuya frontera es una superficie cerrada \(S\) orientada hacia afuera, y \(\mathbf{F}\) es un campo vectorial diferenciable, entonces:

\[
\iint_S \mathbf{F} \cdot \mathbf{n} \, dS = \iiint_V \nabla \cdot \mathbf{F} \, dV
\]

Esta poderosa herramienta permite convertir integrales de superficie complejas en integrales de volumen más sencillas, o viceversa.

🎧 Podcast: 08\_aplicaciones\_teorema-de-la-divergencia-gauss

Tu navegador no soporta el elemento de audio.

\subsection{Ejemplo resuelto: Flujo a través de un elipsoide}\label{ejemplo-resuelto-flujo-a-travuxe9s-de-un-elipsoide}

\begin{quote}
\textbf{Ejemplo 3.} Calcule el flujo del campo \(\mathbf{F}(x, y, z) = (x, y, z)\) a través de la frontera del elipsoide \(\frac{x^2}{a^2} + \frac{y^2}{b^2} + \frac{z^2}{c^2} \leq 1\).
\end{quote}

\textbf{Solución.} La divergencia del campo es:

\[
\nabla \cdot \mathbf{F} = \frac{\partial x}{\partial x} + \frac{\partial y}{\partial y} + \frac{\partial z}{\partial z} = 3
\]

Por el teorema de la divergencia:

\[
\Phi = \iiint_V 3 \, dV = 3 \cdot V_{\text{elipsoide}}
\]

El volumen de un elipsoide de semiejes \(a, b, c\) es \(V = \frac{4}{3}\pi abc\). Por tanto:

\[
\Phi = 3 \cdot \frac{4}{3}\pi abc = 4\pi abc
\]

Notablemente, este resultado generaliza el flujo a través de una esfera (\(a = b = c = r\)), donde \(\Phi = 4\pi r^3\).

\section{Volumen bajo una superficie}\label{volumen-bajo-una-superficie}

Si \(z = f(x, y)\) es una superficie sobre una región \(R\) del plano \(xy\), el volumen del sólido comprendido entre \(R\) y la superficie es:

\[
V = \iint_R f(x, y) \, dx\, dy
\]

\subsection{Ejemplo resuelto}\label{ejemplo-resuelto-26}

\begin{quote}
\textbf{Ejemplo 4.} Calcule el volumen del sólido acotado por el paraboloide elíptico \(z = 4 - x^2 - y^2\) y el plano \(z = 0\).
\end{quote}

\textbf{Solución.} La superficie corta al plano \(xy\) cuando \(z = 0\), es decir, \(x^2 + y^2 = 4\). La región de integración es el disco \(R = \{(x, y) : x^2 + y^2 \leq 4\}\). Entonces:

\[
V = \iint_R (4 - x^2 - y^2) \, dx\, dy
\]

En coordenadas polares:

\[
V = \int_0^{2\pi} \int_0^2 (4 - r^2) \, r\,dr\,d\theta = 2\pi \int_0^2 (4r - r^3) \, dr = 2\pi \left[ 2r^2 - \frac{r^4}{4} \right]_0^2 = 2\pi (8 - 4) = 8\pi
\]

\section{Aplicaciones en física}\label{aplicaciones-en-fuxedsica}

\subsection{Ley de Gauss eléctrica}\label{ley-de-gauss-eluxe9ctrica}

En electrostática, el flujo del campo eléctrico \(\mathbf{E}\) a través de cualquier superficie cerrada \(S\) es proporcional a la carga neta \(Q_{\text{enc}}\) encerrada por dicha superficie:

\[
\iint_S \mathbf{E} \cdot \mathbf{n} \, dS = \frac{Q_{\text{enc}}}{\varepsilon_0}
\]

Esta ley permite calcular campos eléctricos de distribuciones de carga con alta simetría (esférica, cilíndrica, plana).

\subsection{Ley de Ampère-Maxwell}\label{ley-de-ampuxe8re-maxwell}

En magnetismo, la circulación del campo magnético \(\mathbf{B}\) alrededor de una curva cerrada está relacionada con la corriente eléctrica \(I\) y la variación del flujo eléctrico:

\[
\oint_C \mathbf{B} \cdot d\mathbf{l} = \mu_0 I + \mu_0 \varepsilon_0 \frac{d\Phi_E}{dt}
\]

\subsection{Momento de inercia}\label{momento-de-inercia}

El momento de inercia de una lámina superficial \(S\) de densidad \(\rho\) respecto a un eje es:

\[
I = \iint_S \rho \, d^2 \, dS
\]

donde \(d\) es la distancia de cada punto al eje de rotación. Por ejemplo, el momento de inercia de una esfera de radio \(R\) y masa \(M\) respecto a un diámetro es:

\[
I = \frac{2}{5}MR^2
\]

\section{Aplicaciones en arquitectura e ingeniería}\label{aplicaciones-en-arquitectura-e-ingenieruxeda}

\subsection{Estructuras regladas}\label{estructuras-regladas}

Las superficies regladas (hiperboloides de una hoja, paraboloides hiperbólicos, conos) pueden construirse con elementos rectilíneos. Esta propiedad las hace económicas y estructuralmente eficientes:

\begin{itemize}
\tightlist
\item
  \textbf{Torres de enfriamiento}: los hiperboloides de una hoja distribuyen uniformemente las tensiones del viento.
\item
  \textbf{Techos tensados}: los paraboloides hiperbólicos permiten construir cubiertas amplias con poco material.
\end{itemize}

\subsection{Óptica}\label{uxf3ptica}

\begin{itemize}
\tightlist
\item
  \textbf{Espejos parabólicos}: un paraboloide elíptico refleja rayos paralelos a su eje hacia el foco, y viceversa. Esta propiedad se utiliza en telescopios reflectores, antenas satelitales y faros.
\item
  \textbf{Lentes elípticas}: los elipsoides tienen la propiedad de que un rayo que parte de un foco pasa obligatoriamente por el otro foco después de reflejarse.
\end{itemize}

\subsection{Mecánica celeste}\label{mecuxe1nica-celeste}

Según las leyes de Kepler, las órbitas planetarias son elipses con el Sol en uno de los focos. Las superficies equipotenciales gravitacionales de un cuerpo esférico son esferas concéntricas.

\section{Ejercicios propuestos}\label{ejercicios-propuestos-7}

\begin{enumerate}
\def\labelenumi{\arabic{enumi}.}
\tightlist
\item
  Calcule el área de la porción de esfera \(x^2 + y^2 + z^2 = a^2\) que queda dentro del cilindro \(x^2 + y^2 = ax\).
\item
  Halle el flujo del campo \(\mathbf{F} = (x^3, y^3, z^3)\) a través de la esfera de radio \(R\) centrada en el origen.
\item
  Calcule el momento de inercia de una lámina esférica de radio \(R\), masa \(M\) y densidad uniforme, respecto a un diámetro.
\item
  Determine el volumen del sólido acotado por los paraboloides \(z = x^2 + y^2\) y \(z = 18 - x^2 - y^2\).
\item
  Aplique el teorema de la divergencia para calcular el flujo de \(\mathbf{F} = (x^2, y^2, z^2)\) a través de la frontera del cubo \([0, 1] \times [0, 1] \times [0, 1]\).
\end{enumerate}

\chapter{Conclusiones}\label{conclusiones}

🎧 Podcast del capítulo

Tu navegador no soporta el elemento de audio.

A lo largo de este texto se ha presentado un estudio completo de las herramientas matemáticas necesarias para el análisis de superficies en el espacio tridimensional, desde los fundamentos del álgebra vectorial hasta las aplicaciones del cálculo integral y la geometría diferencial.

\section{Resumen de contenidos por capítulo}\label{resumen-de-contenidos-por-capuxedtulo}

\textbf{Capítulo 0 --- Vectores y Geometría Vectorial.} Se establecieron los fundamentos del álgebra lineal en \(\mathbb{R}^3\): vectores libres y coordenados, operaciones vectoriales, producto punto (proyecciones, ángulos, ortogonalidad), producto cruz (áreas, vectores normales) y triple producto escalar (volúmenes).

\textbf{Capítulo 1 --- Rectas y Planos en el Espacio.} Se desarrollaron las ecuaciones vectoriales, paramétricas y simétricas de rectas; las ecuaciones punto-normal y general de planos; fórmulas de distancia (punto a plano, punto a recta, rectas alabeadas); ángulos entre rectas, planos y recta-plano; e intersecciones.

\textbf{Capítulo 2 --- Sistemas de Coordenadas y Parametrización.} Se estudiaron coordenadas polares, cilíndricas y esféricas con sus fórmulas de conversión y jacobianos. Se parametrizaron curvas clásicas en 2D (cicloide, astroide, epicicloide) y 3D (hélice, curva de Viviani), incluyendo el cálculo de longitud de arco.

\textbf{Capítulo 3 --- Funciones de Varias Variables y Derivadas Parciales.} Se definieron funciones de 2 y 3 variables, curvas de nivel, límites y continuidad. Se desarrollaron derivadas parciales con su interpretación geométrica, plano tangente, aproximación lineal, regla de la cadena, derivada direccional y gradiente.

\textbf{Capítulo 4 --- Máximos, Mínimos y Optimización.} Se presentó el criterio del Hessiano para clasificar puntos críticos, extremos absolutos en regiones compactas y el método de multiplicadores de Lagrange para optimización con restricciones.

\textbf{Capítulo 5 --- Superficies Cuádricas.} Se estudiaron exhaustivamente las seis cuádricas principales: elipsoide, hiperboloides (una y dos hojas), paraboloides (elíptico e hiperbólico) y cono elíptico. Cada una con trazas, propiedades, ejemplos resueltos y visualizaciones interactivas.

\textbf{Capítulo 6 --- Integrales Múltiples.} Se desarrollaron integrales dobles y triples, evaluación como integrales iteradas, cambio de orden de integración, cambio de variables (polares, cilíndricas, esféricas) y aplicaciones (áreas, volúmenes, centroides, momentos de inercia).

\textbf{Capítulo 7 --- Geometría Diferencial de Superficies.} Se introdujeron superficies paramétricas, vector normal, plano tangente, primera forma fundamental (métrica), segunda forma fundamental, curvaturas principales, curvatura gaussiana y curvatura media.

\textbf{Capítulo 8 --- Aplicaciones.} Se presentaron integrales de superficie, flujo de campos vectoriales, teorema de la divergencia de Gauss, y aplicaciones en física (ley de Gauss eléctrica, momentos de inercia), arquitectura (estructuras regladas), óptica (espejos parabólicos) y mecánica celeste.

\section{Conclusiones finales}\label{conclusiones-finales}

\begin{enumerate}
\def\labelenumi{\arabic{enumi}.}
\item
  El álgebra vectorial es el lenguaje natural para describir objetos geométricos en el espacio. Los conceptos de producto punto y cruz son fundamentales para definir planos, ángulos y orientaciones.
\item
  Los sistemas de coordenadas alternativos (polares, cilíndricas, esféricas) simplifican enormemente la descripción de objetos con simetría radial o esférica, y sus jacobianos son esenciales para el cambio de variables en integrales.
\item
  Las derivadas parciales y el gradiente permiten analizar cómo cambian las funciones en cada dirección del espacio, siendo el gradiente la herramienta clave para la optimización.
\item
  El Hessiano proporciona un criterio completo para clasificar puntos críticos en funciones de dos variables, mientras que los multiplicadores de Lagrange extienden estas técnicas a problemas con restricciones.
\item
  Las superficies cuádricas constituyen la familia algebraica más importante en \(\mathbb{R}^3\), y su estudio mediante trazas y parametrizaciones establece el puente entre el álgebra y la geometría diferencial.
\item
  El cálculo integral multivariable (dobles y triples) generaliza las nociones de área y volumen, permitiendo calcular magnitudes físicas como masa, centro de masa y momentos de inercia.
\item
  La curvatura gaussiana es un invariante geométrico profundo que clasifica las superficies según su comportamiento local: esférico (\(K > 0\)), hiperbólico (\(K < 0\)) o plano/cilíndrico (\(K = 0\)).
\item
  Las herramientas desarrolladas en este texto tienen aplicaciones directas en ingeniería estructural, diseño óptico, física de campos, gráficos por computadora y mecánica celeste.
\end{enumerate}

\section{Direcciones para profundizar}\label{direcciones-para-profundizar}

\begin{itemize}
\tightlist
\item
  \textbf{Geometría diferencial clásica}: geodésicas, teorema egregio de Gauss, teorema de Gauss-Bonnet.
\item
  \textbf{Superficies mínimas}: catenoides, helicoides, superficies de Scherk.
\item
  \textbf{Topología de superficies}: clasificación de superficies compactas, género, característica de Euler.
\item
  \textbf{Análisis numérico}: métodos de elementos finitos para ecuaciones en derivadas parciales sobre superficies.
\item
  \textbf{Gráficos por computadora}: modelado NURBS, subdivisiones de superficies, mallado de superficies implícitas.
\end{itemize}

\section{Referencias bibliográficas}\label{referencias-bibliogruxe1ficas}

Para profundizar en los temas tratados, se recomiendan los siguientes textos:

\begin{itemize}
\tightlist
\item
  Swokowski, E. W. \emph{Cálculo con Geometría Analítica}. Cengage Learning, 2009.
\item
  Zill, D. G. y Wright, W. S. \emph{Cálculo de Varias Variables}. McGraw-Hill, 2011.
\item
  Lehmann, C. H. \emph{Geometría Analítica}. Limusa, 2012.
\item
  Do Carmo, M. P. \emph{Differential Geometry of Curves and Surfaces}. Prentice Hall, 1976.
\item
  O'Neill, B. \emph{Elementary Differential Geometry}. Academic Press, 2006.
\item
  Stewart, J. \emph{Cálculo Multivariable}. Cengage Learning, 2012.
\end{itemize}

  \bibliography{book.bib,packages.bib}

\end{document}
